\documentclass[11pt,a4paper]{article}
\usepackage[T1]{fontenc}
\usepackage[utf8]{inputenc}
\usepackage{lmodern}
\pdfmapfile{+lm.map}
\pdfmapfile{+cm.map}
\pdfmapfile{+symbols.map}
\usepackage{amsmath,amssymb,amsthm,mathtools}
\usepackage[a4paper,margin=27mm,headheight=14pt]{geometry}
\usepackage{microtype}
\usepackage{xcolor}
\usepackage{enumitem}
\usepackage{bookmark}
\usepackage{fancyhdr}
\usepackage{needspace}
\usepackage{etoolbox}
\usepackage{titlesec}
\usepackage{hyperref}
\hypersetup{colorlinks=true,linkcolor=black,citecolor=black,urlcolor=blue!45!black,pdftitle={Normal dovetails: analytic admission and quantitative constructions},pdfauthor={}}
\pagestyle{fancy}
\fancyhf{}
\fancyhead[L]{\small\scshape Normal dovetails}
\fancyhead[R]{\small\thepage}
\renewcommand{\headrulewidth}{0.3pt}
\fancypagestyle{plain}{\fancyhf{}\fancyfoot[C]{\thepage}\renewcommand{\headrulewidth}{0pt}}
\setlength{\parindent}{0pt}
\setlength{\parskip}{0.5em plus 0.15em minus 0.1em}
\setlength{\emergencystretch}{2em}
\setlist{nosep,leftmargin=1.8em}
\setcounter{secnumdepth}{2}
\setcounter{tocdepth}{2}
\titleformat{\section}{\large\bfseries}{\thesection.}{0.6em}{}
\titleformat{\subsection}{\normalsize\bfseries}{\thesubsection.}{0.6em}{}
\titlespacing*{\section}{0pt}{2.2ex plus 1ex minus .2ex}{0.8ex}
\titlespacing*{\subsection}{0pt}{1.8ex plus .5ex minus .2ex}{0.5ex}
\providecommand{\tightlist}{\setlength{\itemsep}{0pt}\setlength{\parskip}{0pt}}
\DeclareMathOperator{\dist}{dist}
\newcommand{\lt}{<}
\newcommand{\gt}{>}
\allowdisplaybreaks[2]
\widowpenalty=10000
\clubpenalty=10000
\title{\vspace{-1.5em}\LARGE\bfseries Normal dovetails\par\vspace{0.35em}\large Analytic admission and quantitative constructions}
\author{}
\date{}
\begin{document}
\maketitle
\vspace{-3em}
\begin{abstract}

We concatenate the complete first $A_rD_r$ fractional radix digits of an analytic function at all affine inputs $(bj+1)b^{-A_r}$ in each stage $b^{r-1}\le j< b^r$. An exponential upper bound and a finite-window, factorial-scale lower bound on Taylor coefficients suffice for normality through every output prefix. A single computable schedule works pointwise for all nonpolynomial real D-finite germs on a common positive-radius disk. The schedule can meet any specified computable, nondecreasing, divergent overhead budgets above logarithmic input precision and copied length. Uniformity of the discrepancy estimate requires uniform coefficient-window data; uniformity of the schedule alone does not provide it.

The construction admits rational, algebraic and transcendental analytic source functions. For $(1+x)^\alpha$ the exact exceptional parameters are the nonnegative integers; for $(1-x)^{-\theta}$ they are the nonpositive integers. These are results about newly constructed digit streams, not a claim that normality implies transcendence or a normality theorem for the sampled function values or their limit.

The coefficient-window theorem is the broadest analytic admission result in this account, so we state it first. Singular and flat functions need different derivative estimates and schedules; they are treated afterward rather than being silently included in its hypotheses.

\end{abstract}
\vspace{1em}
\tableofcontents
\newpage

\section*{Introduction}

A real number is normal to base \(b\) if each word \(w\) of length \(W\)
occurs in its fractional expansion with limiting frequency \(b^{-W}\).
Equivalently, its radix orbit is uniformly distributed modulo one. We
ask when many short observations of a function, assembled as consecutive
digit blocks, yield a normal number. The observations here use inputs
approaching zero; each block contains a complete initial segment of the
fractional expansion, and no row is omitted.

The central mechanism is cancellation in the phases at every actual
digit offset. A finite window of Taylor coefficients prevents all the
relevant high derivatives from becoming too small. An order-uniform
derivative estimate then controls exponential sums, and a Fourier
approximation converts cancellation into word counts. Estimates for
incomplete row intervals are what extend the conclusion from stage
endpoints to every output prefix.

This model differs from the integer-part concatenations studied by
Madritsch, Thuswaldner and Tichy \cite{madritsch}, where an entire
function is sampled at natural numbers or primes tending to infinity. We
use the growing-order derivative estimate of Arias de Reyna
\cite{arias}. Normality and transcendence remain separate questions: the
classical fixed-base normal Liouville construction in Appendix C is
included as background, while Becher, Heiber and Slaman \cite{becher}
construct a computable absolutely normal Liouville number.

Theorem A gives the general analytic criterion. Sections 3 and 4 derive
D-finite and exact power-family consequences. Sections 5--7 treat
singular sources, two flat kernels, and the scale of approach to
polynomial exceptions. Section 8 compares different quantitative
objectives. The appendices supply the derivative estimate, endpoint-safe
Fourier transfer, and elementary transfer and obstruction arguments.

\section{Construction and the coefficient-window
theorem}\label{construction-and-the-coefficient-window-theorem}

Fix an integer \(b\ge2\). Stage \(r\ge1\) contains every row index
\(b^{r-1}\le j< b^r\). Choose positive integers \(A_r,D_r\) and put
\(K_r=A_rD_r\). For a row \(j\) in stage \(r\), use the aliases
\(A_j=A_r\), \(D_j=D_r\), and \(K_j=K_r\). The source input is

\[
x_j=(bj+1)b^{-A_r},\qquad b^{r-1}\le j< b^r.
\]

For each fixed parameter \(u\) in a set \(U\), reduce the real source
value \(h_u(x_j)\) modulo 1 and copy its complete first \(K_r\)
fractional digits. Use the terminating expansion at ambiguous endpoints.
Concatenating these blocks in increasing \(j\) gives a number \(Y_u\).
For a word \(w\) of length \(W\), let \(C_w(N;u)\) count occurrences
starting at the first \(N\) positions of this infinite output, including
words that end beyond position \(N\). Counting only words completed by
\(N\) changes the count by at most \(W-1\). The parameter is fixed
within each output; it never varies from row to row.

An \textbf{admissible schedule} retains all rows in their original
order, uses positive integers \(A_r,D_r\), copies each full prefix of
\(K_r\) digits, and satisfies \(x_j\to0\) and \(D_r\to\infty\). Every
sampled value must lie in the stated real domain of the source.
Additional budget or regularity conditions are imposed only where
explicitly stated.

A \textbf{common schedule} does not depend on \(u\). \textbf{Pointwise
normality} means that every fixed parameter gives a normal output.
\textbf{Uniform normality on \(U\)} means that the normalized word-count
error tends to zero uniformly over \(u\in U\); compact-uniform normality
requires this on each specified compact parameter set. A common schedule
need not give uniform error near exceptional parameters. A computable
schedule also does not by itself make the copied digits computable from
arbitrary real input data.

In applications, a parameterized source may be written as
\(h_\theta(x)=\mathcal B(\theta,g(x))\). The operation \(\mathcal B\)
acts before digits are copied; it is not applied to an already normal
output.

\subsection{Coefficient-window
hypotheses}\label{coefficient-window-hypotheses}

A family of real analytic germs \(h_u\) on a common complex disk
\(|x|< \rho\) around zero is \textbf{window admissible on a parameter
set \(U\)} if, writing

\(h_u'(x)=\sum_{n\ge 0} c_n(u) x^n\), it admits a specified positive
rational lower bound \(\rho\) for the common analytic radius, numbers
\(M\ge 1\), \(C\ge 1\), integers \(d\ge 1\), \(n_0\ge 0\), independent
of \(u \in U\), such that, with the shared exponential scale
\(S=\max(2,2/\rho)\):

(W1) \(|c_n(u)| \le  M S^n\) for all \(n\ge 0\) and all \(u \in U\);

(W2) for every \(n\ge n_0\),

\[
\max_{0\le i< d} |c_{n+i}(u)| \ge  \exp(-C(n+d+1) \log(n+d+2))
\]

The coefficients and the germs are real, and their Taylor series
represent the actual source functions on this disk. Cauchy's estimate on
the one interior disk \(|x|\le \rho/2\) gives this shared exponential
scale \(S\); a uniform bound there supplies \(M\) in (W1). Pointwise,
\(M\) may depend on the germ. The other window constants may be merely
existential for qualitative statements, or computably supplied for an
effective explicit onset and compact-uniform bound. No rationality
hypothesis is imposed on coefficients or parameters.

\Needspace{6\baselineskip}
\textbf{Theorem A (window admission, strong all-prefix version).} Let
\(h_u\) satisfy (W1)-(W2) on \(U\) with the specified common
analytic-radius lower bound \(\rho\). Let
\(F,G:\mathbb N_{>0}\longrightarrow\mathbb N_{>0}\) be computable
nondecreasing functions tending to infinity. There is one computable
schedule, depending only on \(b,F,G,\rho\) and \textbf{not on \(u\) or
the remaining window data}, with:

\[
x_j\longrightarrow0,\qquad D_j\longrightarrow\infty,\qquad \frac{A_j}{\log_bj}\longrightarrow1,
\]

and, eventually,

\[
A_j-\log_bj\le G(j),\qquad K_j\le F(j)\log_bj.
\]

such that, for every fixed word \(w\), uniformly in \(u \in U\),

\[
\left|C_w(N;u)-Nb^{-W}\right|\le C_{U,b,W}N\frac{B_{r(N)}}{r(N)D_{r(N)}}
\]

for all sufficiently large \(N\). Here \(r(N)\) is the stage containing
the last observed output digit and \(B_r=A_r-r\), as defined precisely
below. In particular the concatenation is normal uniformly on \(U\). For
any shared positive domain radius, the same schedule works
\emph{pointwise} for every germ having some finite data satisfying
(W1)-(W2) on that common domain, even when the coefficient-window
constants cannot be chosen uniformly across the germs. Without a fixed
positive radius lower bound, a domain-dependent finite initial
adjustment is required for each germ: a fixed early sample may be
outside its domain. This adjustment leaves all eventual budgets and
discrepancy bounds unchanged.

The theorem holds equally for \(h_u(x)\) with an additive integration
constant or an additional additive center, because these only rotate
each harmonic sum at a common digit offset. That observation is a proof
convenience, not a replacement of \(\mathcal B(u,g(x))\) by an additive
center.

\section{Proof of the coefficient-window
theorem}\label{proof-of-the-coefficient-window-theorem}

The arithmetic of a rational Taylor series can support fresh-prime
grids, but it is not the only way to obtain cancellation. Here the phase
is differentiated without truncating the source. The order of
differentiation grows, so a derivative estimate with constants uniform
in that order is essential.

A lower bound on every coefficient would exclude functions such as
cosine. A lower bound on the largest coefficient in each fixed-width
window allows zeros and cancellations. An invertible
polynomial-coefficient recurrence supplies precisely this weaker
condition by running the recurrence backwards. This is the mechanism
behind the D-finite corollary.

Here ``rational'', ``algebraic'' and ``transcendental'' describe source
functions, not their output numbers. Whether any algebraic irrational is
normal remains open, and no implication from normality to transcendence
is used.

\subsection{Explicit schedule without a prescribed convergence
modulus}\label{explicit-schedule-without-a-prescribed-convergence-modulus}

Set

\[
S=\max\left(2,\frac2\rho\right),\qquad B_0=3+\left\lceil\log_b(4S)\right\rceil.
\]

For \(r\ge2\), define

\[
B_r=\max\left\lbrace B_0,\,
\min\left(\left\lfloor\sqrt r\right\rfloor,
\left\lfloor\frac{G(b^{r-1})}{4}\right\rfloor\right)\right\rbrace,
\]

\[
\begin{aligned}
D_1&=1,\\
D_r&=\min\left\lbrace D_{r-1}+1,
\max\left(1,\left\lfloor\frac{\log_2B_r}{8}\right\rfloor\right),
\max\left(1,\left\lfloor\frac{F(b^{r-1})}{4}\right\rfloor\right)
\right\rbrace,\\
A_r&=r+B_r,\qquad K_r=A_rD_r.
\end{aligned}
\]

All floors, ceilings and thresholds can be evaluated by integer or
rational comparisons. No rate of divergence for \(F\) or \(G\) is
needed.

For stage 1 take \(B_1=B_0, A_1=1+B_0, D_1=1\) (only a finite
adjustment). This ensures \(x_j\le b^{2-B_0}< 1/(2S)\) from the first
row. These are nondecreasing, divergent \(B\) and \(D\) profiles,
\(B_r\le \sqrt{r}\) for all sufficiently large \(r\),
\(D_r\le 1+(\log_2 B_r)/8\), and \(2^D=O(B^{1/8})\). The
min-with-previous-plus-one guarantees \(D_r\le D_{r-1}+1\) and no
downward jumps. For \(j\) in stage \(r\), \(\log_b j\) lies in
\([r-1,r)\), so the two eventual budgets follow from the caps. Since
\(x_j\le b^{2-B_r}\), the source lies in the domain and tends to zero.
Since \(K_r=(r+B_r)D_r\), both \(A_r/\log_b j \longrightarrow  1\) and
\(K_r/\log_b j \longrightarrow  \infty\) hold, and
\(B_r/(rD_r)\longrightarrow 0\).

\subsection{Uniform analytic derivative
lemma}\label{uniform-analytic-derivative-lemma}

For \(k\ge 1\), let \(c_{k-1}\) be a usable coefficient furnished by
(W2). For \(0\le x< 1/S\), the Taylor remainder identity gives

\[
\frac{h_u^{(k)}(x)}{(k-1)!}=c_{k-1}+\sum_{\ell\ge1}\binom{k+\ell-1}{\ell}c_{k+\ell-1}x^\ell
\]

Therefore, if \(2k S x\le 1\),

\[
\left|\frac{h_u^{(k)}(x)}{(k-1)!}-c_{k-1}\right|\le4MkS^kx \tag{2.1}
\]

This follows from the majorant \(M S^{k-1}[(1-Sx)^{-k}-1]\),
\(-\log(1-z)\le 2z\), and \(e^v-1\le 2v\) for \(0\le v\le 1\). Condition
(W2) converts the absolute error to a relative error at most \(1/2\)
whenever

\[
B_r\ge B_*(M,S,C,d,n_0,b)+4Ck\log(k+2)
\]

Only \(k\le D_r+n_0+2d+12=O(\log B_r)\) is needed, so this inequality
holds for all sufficiently large \(r\), uniformly on \(U\). Consequently
\(h_u^{(k)}\) has one fixed sign throughout every required real row
interval, and lies between one half and three halves of
\((k-1)! c_{k-1}\) in signed magnitude. Furthermore, for every selected
order,

\[
\left|\log_b\bigl((k-1)!\,|c_{k-1}|\bigr)\right|\le\frac{B_r}{16} \tag{2.2}
\]

for all sufficiently large \(r\). When effective upper and nonzero lower
bounds for the window data are supplied, these thresholds can be found
by finite integer/rational bounds. With merely existential data the
thresholds are existential; the common schedule is still computable from
\(b,F,G,\rho\). No hidden order-dependent van der Corput constant is
used.

\subsection{Actual-phase Fourier
calculation}\label{actual-phase-fourier-calculation}

Fix a late stage, abbreviate \(B=B_r\), \(A=r+B\), \(D=D_r\), \(K=AD\),
\(m_{\mathrm{full}}=(b-1)b^{r-1}\); include a real interval containing
\(b^{r-1}-1 \le  y \le  b^r\). For the word length \(W\), discard the
first \(T B\) offsets of each copied row and the final \(W-1\) crossing
offsets, with

\(L_0=n_0+d+10\), \(L_1=L_0+d+1\), \(T=8(L_1+1)\).

For each retained integer offset \(TB\le t\le K-W\), set
\(z=\left\lfloor t/(A-1)\right\rfloor\). Condition (W2) supplies some
index \(n\) in \([z+L_0,z+L_0+d-1]\) with \(c_n \ne  0\) and the stated
lower bound; take \(k=n+1\). Then

\(3\le k\le D+L_1\), \(t-k(A-1)\le -L_0(A-1)\),

\(t-kB \ge  t(1-B/r)-L_1 B \ge  5t/8\) for \(B/r\le 1/4\). (2.3)

For harmonic \(h\ge 1\), define the \textbf{untruncated actual phase}

\[
f_{t,h}(y)=h b^t h_u((by+1)b^{-A})
\]

(The letter h in the subscript is a harmonic integer, not the source
germ.) The chain rule and (2.1)-(2.3) give derivatives of one sign and
\(\lambda_h \le  |f^{(k)}| \le  \Lambda_h\) with
\(\Lambda_h/\lambda_h=3\). For \(h=1\) and all sufficiently large
stages,

\(\lambda_1 m_{\mathrm{full}}^k \ge  b^{t/2}\),
\(\Lambda_1 \le  b^{-8r}\). (2.4)

The first inequality uses \((k-1)!|c_{k-1}| \ge  b^{-B/16}\) and
\(m_{\mathrm{full}}\ge b^{r-1}\). The second uses \(k\ge z+L_0+1\),
\(L_0\ge 11\), and the upper side of (2.2). No value of \(u\) enters
(2.3)-(2.4) except through constants already uniform in (W1)-(W2).

Put \(Q_r=2^{D+L_1}\) and \(s_t=\min(t,r)\). Choose a \textbf{separate
integer Fejér cutoff at every actual offset}:

\[
H_t=\max(1,\left\lfloor b^{s_t/(8 Q_r)}\right\rfloor)
\]

For \(1\le h\le H_t\), \(\Lambda_h\le b^{-7r}\),
\(\lambda_h m_{\mathrm{full}}^k\ge b^{t/2}\), and
\(\Lambda_h/\lambda_h=3\). Applying the growing-order derivative
estimate in Appendix A to any consecutive row subinterval with more than
k terms, and evaluating its absolute upper bound at
\(m_{\mathrm{full}}\), yields

\[
|\sum_{j \in I} \exp(2 \pi i h b^t h_u(x_j))| \le  12 m_{\mathrm{full}} b^{-s_t/Q_r} \tag{2.5}
\]

Indeed the three derivative expressions are bounded by \(b^{-r/2^k}\),
\(b^{-6r/2^k}\), and \(b^{-t/2^k}\). The bound holds also for intervals
with at most \(k\) terms by the trivial estimate
\(k\le 12 m_{\mathrm{full}} b^{-s_t/Q_r}\), valid uniformly after
another effective finite stage threshold. Importantly, (2.5) applies at
each \textbf{actual} offset, to the \textbf{complete untruncated phase},
for all required harmonics and every incomplete consecutive row
interval. No selected-row reindexing or independent digit approximation
is made.

\subsection{Offset-dependent Fourier-to-words and
summation}\label{offset-dependent-fourier-to-words-and-summation}

Use the endpoint-safe Fejér transfer proved in Appendix B: for \(m\)
circle phases \(z_1,\ldots,z_m\) and any interval of length \(b^{-W}\),
the word error is at most

\[
\frac{2m}{\sqrt{H+1}}+\frac2\pi\sum_{h=1}^H\frac{\left|\sum_j e(hz_j)\right|}{h}
\]

A word beginning at offset \(t\) within a row is exactly membership of
\(\lbrace b^t h_u(x_j)\rbrace\) in its radix cylinder in base \(b\); no
digit-near-boundary implication is used. From (2.5) and the chosen
\(H_t\), the retained-offset word contribution is

\[
\le  2m b^{-s_t/(16 Q_r)} + 8 m_{\mathrm{full}} b^{-s_t/Q_r}(1+(\log b)s_t/(8Q_r)) \tag{2.6}
\]

For fixed \(b\ge 2\), geometric summation over \emph{all} offsets
\(TB\le t\le K-W\) gives

\[
\sum_t b^{-s_t/(16Q_r)} \le  C_b Q_r + K b^{-r/(16Q_r)}
\]

\[
\sum_t b^{-s_t/Q_r}(1+s_t/Q_r) \le  C_b Q_r + K(1+r/Q_r)b^{-r/Q_r}
\]

Because \(Q_r=O(B_r^{1/8})=O(r^{1/16})\) and \(K_r=O(r \log r)\), each
final tail is \(O(Q_r)\) after an effective stage threshold. Thus the
total bound on any \textbf{nonempty} consecutive subbatch of \(m\) rows
(in particular the original initial segment of that stage) is

\[
\text{error}\le m(TB_r+W+C_bQ_r)+m_{\mathrm{full}}C_bQ_r \tag{2.7}
\]

An empty subbatch contributes zero. A single incomplete final copied row
is charged its actual number \(z\le K_r\) of observed starts. All
crossing starts are charged inside the \(W\) term. Equation (2.7),
rather than a batch-endpoint-only assertion, is the source of the
all-prefix theorem.

\subsection{Every output prefix}\label{every-output-prefix}

For \(N\) whose last digit lies in stage \(r\), sum (2.7) over all
completed stages and the completed portion of stage \(r\), with only
trivial \(m_s K_s\) charges for the finitely many early stages preceding
the analytic onset. Because \(B_s\) and \(D_s\) are nondecreasing,
\(Q_s\le Q_r\), and the full row counts grow geometrically, this gives

\[
\left|C_w(N;u)-Nb^{-W}\right|=O_{U,b,W}\!\left(b^r(B_r+Q_r+W)+K_r\right) \tag{2.8}
\]

For \(r\ge 3\), the preceding full stage is already part of the prefix
and gives

\[
N \ge  (b-1)b^{r-2} K_{r-1} \ge  c_b b^r r D_r
\]

using \(D_r\le D_{r-1}+1\le 2 D_{r-1}\) and \(K_{r-1}\ge (r-1)D_{r-1}\).
Since \(K_r=O(rD_r)\) and \(Q_r=O(B_r^{1/8})=O(B_r)\), dividing (2.8) by
\(N\) proves Theorem A with \(O(N B_r/(rD_r))\). This argument includes
all original source rows, all word lengths fixed in advance, any
incomplete last stage, and any incomplete last row. The early-stage
contribution is finite and uniform on \(U\) when its admission constants
are uniform.

For example, let \(L_r\) be the least integer with \(2^{L_r} \ge  r+2\)
and use the unbudgeted exact-integer choice
\(B_r=\max(B_0,\min\lbrace m\ge 1: m^3\ge L_r\rbrace )\) and the
\textbf{maximal bounded-increment} depth recurrence

\(D_1=1\),
\(D_r=\min(D_{r-1}+1,\max(1,\left\lfloor (\log_2 B_r)/8\right\rfloor))\),

with no additional smaller F cap. This gives \(D_r=\Theta(\log B_r)\)
eventually and therefore the quantitative uniform rate

\[
O_{U,b,W}\!\left(\frac{N(\log\log N)^{1/3}}{\log N\,\log\log\log N}\right)
\]

here all iterated logarithms are used only in the sufficiently large
\(r\)egime. An unspecified increasing depth merely lying below that cap
need not have this rate. Using
\(B_r=\max(B_0,\min\lbrace m\ge 1: m^2\ge L_r\rbrace )\) with the same
maximal recurrence gives the weaker but cleaner
\(O\!\left(\frac{N\sqrt{\log\log N}}{\log N}\right)\). The
arbitrary-budget rule above can make B diverge even more slowly, with an
exact profile-dependent envelope.

\textbf{Scope caution.} This theorem is proved for functions meeting
(W1)-(W2). It is not a blanket theorem for all analytic functions:
deliberately sparse or super-small coefficient sequences can violate the
window estimate.

\section{D-finite functions and the rational-derivative
classification}\label{d-finite-functions-and-the-rational-derivative-classification}

\Needspace{6\baselineskip}
\textbf{Lemma B (quantitative recurrence window).} Suppose, for every
\(n\ge n_0\), the sequence of Taylor coefficients satisfies

\[
q_d(n,u)c_{n+d}(u)+ \ldots + q_0(n,u)c_n(u)=0
\]

where both endpoint coefficients \(q_0(n,u)\) and \(q_d(n,u)\) are
nonzero for every indicated \(n\) and \(u\). Suppose \(d,n_0\) are fixed
on \(U\), and for constants \(P\ge 1\), \(s\ge 0\) shared on \(U\),

\[
\sum_{i=1}^d |q_i(n,u)/q_0(n,u)| \le  P(n+2)^s
\]

and the initial window vector
\(V_{n_0}(u)=(c_{n_0}(u),\ldots,c_{n_0+d-1}(u))\) has supremum norm at
least \(\nu>0\) uniformly on \(U\). Then for \(n\ge n_0\),

\[
\max_{0\le i< d}|c_{n+i}(u)| \ge  \nu P^{-(n-n_0)} ((n+1)!)^{-s} \tag{3.1}
\]

and hence (W2) holds with a uniform constant \(C\) depending on
\(\nu,P,s,n_0,d\), computable when effective upper bounds for \(P,s\)
and a positive lower bound for \(\nu\) are supplied.

\textbf{Proof.} Let \(V_n=(c_n,\ldots,c_{n+d-1})\). The companion step
\(V_{n+1}=T_n V_n\) is invertible because both endpoint recurrence
coefficients are nonzero. Its inverse shifts all but the recovered first
coordinate; the first coordinate is \(-\sum_{i=1}^d(q_i/q_0)c_{n+i}\).
Hence
\(\lVert T_n^{-1}\rVert_\infty\le\max\left(1,\sum_{i=1}^d|q_i/q_0|\right)\le P(n+2)^s\).
Iterating \textbf{backwards} from \(V_n\),

\[
\nu\le\lVert V_{n_0}\rVert_\infty\le\left(\prod_{t=n_0}^{n-1}P(t+2)^s\right)\lVert V_n\rVert_\infty
\]

which gives (3.1) by bounding the product by \(((n+1)!)^s\). The
factorial logarithm is at most \((n+1)\log(n+1)\), giving (W2). This
proof bounds the window norm, not each coefficient; cancellations in
individual coefficients are allowed. \(\square\)

\Needspace{6\baselineskip}
\textbf{Corollary C (real analytic D-finite derivatives).} Let \(h\) be
real analytic on a neighborhood of zero and suppose its derivative
\(h'\) is D-finite over \(\mathbb R(x)\) (equivalently \(h\) itself is
D-finite), and \(h'\) is not a polynomial. Then \(h\) meets (W1)-(W2)
for some finite constants. In particular Theorem A gives a computable
schedule for every such \(h\) once a positive rational analytic-radius
lower bound is fixed; that schedule needs no recurrence computation. On
a class with fixed positive common analytic radius, a single schedule
can be selected \textbf{without knowing its D-finite recurrence} and
works pointwise for that entire class.

\textbf{Proof.} A linear differential equation with polynomial
coefficients for \(h'\) induces a finite-width recurrence with
polynomial coefficients in \(n\), valid for all sufficiently large
\(n\). Remove identically zero outer coefficient polynomials and shift
the recurrence so that its first and last coefficients are not
identically zero. Each has only finitely many integer roots, so past a
finite \(n_0\) both are nonzero. For arbitrary real recurrence
coefficients this cutoff is existential. Computing it requires effective
coefficient bounds and certified nonzero lower bounds for the leading
nonzero coefficients (or equivalent effective root-bound data); a
symbolic real recurrence alone does not supply those. Effective window
constants also require effective initial-vector nondegeneracy and
analytic upper bounds. Ratios \(q_i/q_0\) are at most a constant times
\((n+2)^s\), because these are real polynomials in \(n\), and the
initial window cannot be identically zero: otherwise the nonzero leading
\(q_d\) forces all later coefficients to vanish and \(h'\) would be
polynomial. An analytic radius provides (W1) by Cauchy's estimate. For
compact parameter families whose recurrence coefficients, radius and
initial vector obey shared nondegeneracy and bound conditions, this
argument supplies the same constants uniformly. On any class with a
specified common analytic disk, the theorem's schedule uses only
\(b,F,G,\rho\), through \(S=\max(2,2/\rho)\), so it applies pointwise
even if the remaining coefficient-window constants vary arbitrarily
between individual sources. Without a common disk, individual schedules
may be chosen; a single fixed finite first input need not lie in every
germ's domain. \(\square\)

The D-finite hypothesis alone does \textbf{not} provide uniformity
across a degenerating parameter family: the first/last recurrence
coefficients or the initial window may tend to zero, and a common
coefficient lower bound can fail. Composition of arbitrary D-finite
functions is not asserted to be D-finite. Products of D-finite germs and
algebraic substitutions in suitable analytic domains furnish standard
structured examples; explicit recurrences should be checked rather than
assumed.

Examples covered by Corollary C include analytic primitives of \(e^x\),
\(\cos x\), \(\log^2(1+x)\) and their nonpolynomial real linear
combinations (if a combination becomes polynomial, it leaves the
admitted class). For \(\cos x\), every other Taylor coefficient
vanishes, but the recurrence \(c_{n+2}=-c_n/((n+1)(n+2))\) has an
invertible two-dimensional window.

\subsection{Exact rational-derivative
admission}\label{exact-rational-derivative-admission}

\Needspace{6\baselineskip}
\textbf{Corollary D (rational-derivative admission).} Let
\(R\in\mathbb Q(x)\) be regular at zero, and let \(g(0)=0\), \(g'=R\).
In the complete-prefix affine-input model, some schedule normalizes
\(\theta+g(x_j)\) uniformly over all additive centers \(\theta\) if and
only if \(R\) is not a polynomial. In the positive case the schedule in
Theorem A meets any chosen budgets \(F,G\) there.

\textbf{Proof.} A nonpolynomial rational function is D-finite and
analytic on a disk about zero; so is its primitive. Corollary C and
Theorem A apply. Additive centers do not change the derivative
coefficients and multiply each harmonic sum at a common offset by a
complex number of modulus one. Thus the discrepancy estimate is uniform
in the center. A positive rational radius bound is effective from the
rational numerator and denominator after cancellation: if
\(Q(x)=q_0+\sum_{i\ge1}q_ix^i\) with \(q_0\ne0\), choose rational
\(0< \rho\le1\) with \(\rho\sum_{i\ge1}|q_i|< |q_0|\). Then \(Q\) has no
zero on \(|x|\le\rho\).

Conversely, if \(R\) is polynomial, \(g\) is a rational-coefficient
polynomial of some degree \(m\) (including the zero polynomial). Choose
a positive integer \(q\) clearing its coefficients and set the center to
zero. For every row in stage \(r\) and integer offset \(t\ge mA_r\), the
number \(q b^t g(x_j)\) is an integer. Every deep source phase therefore
belongs to the fixed grid \(q^{-1}\mathbb Z/\mathbb Z\). Choose a fixed
word length \(W\) with \(b^W>q\). Some canonical radix cylinder of
length \(W\) misses that grid, and that word never starts at any deep
within-row position. The fraction of remaining positions at
completed-stage endpoints tends to zero: each row loses at most
\(mA_r+W\) starts out of \(A_rD_r\), and \(D_r\to\infty\). The finitely
many early rows are negligible. The omitted word has frequency zero
rather than \(b^{-W}\), contradicting normality. Thus center-uniform
normality is impossible. \(\square\)

This is a function-class classification for existence of a schedule, not
a classification of all schedules. The primitive \(g(x)=x/(1-x)\) is
admitted, despite being rational. At the same time, the inverse
operation \(g/(1+g)=x\) destroys admission at center zero. Arbitrary
nonlinear composition therefore does not preserve the property.

\section{Parameterized operations and exact power
exceptions}\label{parameterized-operations-and-exact-power-exceptions}

For a specified parameter set \(U\), call \(\mathcal B\)
\textbf{normalization-surviving for g} if some admissible common
schedule makes every \(Y_u\), \(u \in U\), normal; call it
\textbf{compact-uniformly surviving} if each prescribed compact subset
has a uniform \(o(N)\) bound under that same schedule. A parameter is a
\textbf{structural exception} if no admissible schedule normalizes its
source. \textbf{Schedule suppression} means a particular schedule
produces a nonnormal output at the parameter, without asserting
structural failure. \textbf{Exact cancellation} means an
algebraic/functional identity eliminates high derivatives, while
\textbf{near-cancellation} refers to a quantitative, observation-scale
loss; the last three notions are not interchangeable.

\Needspace{6\baselineskip}
\textbf{Theorem E (two regular powers).} In the observation model of
Section 1 the following classifications hold, each with a single
computable schedule independent of its parameter, normal pointwise for
every listed nonexceptional real parameter and uniform on compact sets
disjoint from the listed exceptions:

\begin{enumerate}
\def\labelenumi{(\roman{enumi})}
\item
  \(\mathcal B(\theta,y)=\theta^y\), \(\theta>0\), \(g(x)=\log(1+x)\):
  write \(\alpha=\log \theta\), take the positive real branch;
  \(h_\theta(x)=(1+x)^\alpha\). The exact exception set is
  \(\theta=e^m\), \(m=0,1,2,\ldots\).
\item
  \(\mathcal B(\theta,y)=(1+y)^\theta\), \(g(x)=x/(1-x)\) on
  \(0< x< 1\): \(h_\theta(x)=(1-x)^{-\theta}\). The exact exception set
  is \(\theta=0,-1,-2,\ldots\).
\end{enumerate}

\textbf{Proof.} For (i), \(h'(x)=\sum c_n x^n\) with

\(c_0=\alpha\), \(c_{n+1}/c_n=(\alpha-1-n)/(n+1)\).

For \(\alpha\) not a nonnegative integer, \(c_0\ne 0\) and no recurrence
factor vanishes. The one-dimensional backward recurrence on any compact
parameter set avoiding these integers supplies (W2); a common disk
\(|x|< 1/2\) supplies (W1). For \(\alpha=m\ge 0\), \((1+x)^m\) is an
integer-coefficient polynomial. In every row, at offsets \(t\ge m A_r\),
\(b^t h(x_j)\) is an integer; therefore all but \(m A_r+O(W)\) of its
\(K_r=A_rD_r\) digit starts are strings of zeros. Since
\(D_r\longrightarrow \infty\), a fixed all-zero word has frequency
tending to 1 at completed-stage endpoints rather than \(b^{-W}\).

For (ii), \(h'(x)=\theta(1-x)^{-\theta-1}\), with

\(c_0=\theta\), \(c_{n+1}/c_n=(\theta+n+1)/(n+1)\).

The same argument applies; at \(\theta=-m\) the value \((1-x)^m\) is an
integer-coefficient polynomial, with exactly the same fixed-center
obstruction. The classification is exact in the sense of existence of a
schedule in this model: these polynomial endpoints fail for
\textbf{every} admissible schedule, not just the theorem's one.
\(\square\)

The classification concerns these explicit compositions, not an
arbitrary complex or negative-base real power. Parameter \(\theta\) is
held fixed along each output.

\section{Singular powers on positive
inputs}\label{singular-powers-on-positive-inputs}

Consider \(\mathcal B(\theta,y)=y^\theta\), \(g(x)=x/(1-x)\) and
\(x>0\). Choose the positive real branch
\(y^\theta=\exp(\theta \log y)\) for every real \(\theta\). Then

\[
h_\theta(x)=x^\theta (1-x)^{-\theta}.
\]

\Needspace{6\baselineskip}
\textbf{Theorem F (singular powers).} There is one computable stage
schedule of the form in Theorem A for which every fixed
\textbf{\(\theta\ne 0\)} produces a normal concatenation, all original
rows and cutoffs included. For compact \(J\) avoiding the nonnegative
integers, the discrepancy estimate of Theorem A is uniform on \(J\). The
positive integer parameters are also admitted by the analytic
rational-derivative theorem, using the \textbf{same source schedule},
but a bound uniform while crossing those integers is not established
here. At \(\theta=0\) the source equals 1 and every fractional digit is
zero.

\textbf{Proof on compact \(J\) disjoint from nonnegative integers.} Set
\(u_\theta(x)=(1-x)^{-\theta}\) and
\(p_k(\theta)=\theta(\theta-1)\cdots(\theta-k+1)\). The quantity

\[
\delta=\min_{\substack{\theta\in J\\ \ell\in\mathbb Z_{\ge 0}}}|\theta-\ell|
\]

is strictly positive. Leibniz's formula gives, for \(k\ge 1\),

\[
h_\theta^{(k)}(x)=p_k(\theta) x^{\theta-k} u_\theta(x) [1+E_{k,\theta,x}],
\]

where

\[
|E_{k,\theta,x}| \le  \exp(C_J k^2 x/\delta)-1.\tag{5.1}
\]

To check (5.1), divide the Leibniz term of order \(\ell\) by the
\(\ell=0\) term: the ratio of falling factorials contributes at most
\(\delta^{-\ell}\), while
\(u_\theta^{(\ell)}/u_\theta=(\theta)^{\overline{\ell}}(1-x)^{-\ell}\)
contributes at most \((\ell+C_J)^\ell(1-x)^{-\ell}\). Together with
\(\binom{k}{\ell}\le k^\ell/\ell!\) the sum is bounded by an exponential
in \(C_J k^2 x/\delta\). Since \(x\le b^{2-B_r}\) and
\(k\le D_r+L_1=O(\log B_r)\), it is at most 1/2 for large \(r\),
uniformly on \(J\). Thus the derivative has one sign. The lower product
satisfies \(|p_k(\theta)|\ge \min(1,\delta)^k\); the upper satisfies
\(|p_k(\theta)|\le (k+C_J)^k\).

On the real source interval \(b^{r-1}-1\le y\le b^r\), the input
\(x=(by+1)b^{-A_r}\) lies between \(b^{-B_r-1}\) and \(b^{2-B_r}\) for
large \(r\). Consequently the variation ratio of the derivative of order
\(k\) is at most \(6 b^{3(k+C_J)}\). This is \textbf{not} a constant 3,
but \(\log_b\) of it is \(O_J(k)=o(r)\). Taking
\(k=\left\lfloor t/(A_r-1)\right\rfloor+L_1\) with a sufficiently large
fixed \(L_1=L_1(J)\) gives \(3\le k\le D_r+L_1\) for all sufficiently
large \(r\) and all retained \(t\ge T_J B_r\). The floor gives precisely

\[
t-k r \le  -(L_1-1)r+O(D_r B_r),
\]

because \(\left\lfloor t/(A_r-1)\right\rfloor\ge t/(A_r-1)-1\); the
stronger \(-L_1 r\) is not justified. A sufficiently large fixed \(L_1\)
absorbs this one-unit loss, and \(D_r B_r=o(r)\). First apply the
derivative estimates to harmonic 1:

\[
\lambda_1 m_{\mathrm{full}}^k \ge  b^{t/2},\qquad \Lambda_1\le b^{-8r},\qquad
\frac{\Lambda_1}{\lambda_1} \le  6b^{3(k+C_J)}.\tag{5.2}
\]

The lower estimate reflects cancellation of \(x^{-k}\) against the
\(b^{-k B}\) factor:
\(\log_b(\lambda_1 m_{\mathrm{full}}^k) = t - \theta B_r + O_J(k \log(k+2))\),
uniformly on the interval. The upper follows from the floor inequality
and a sufficiently large \(L_1\). Choose \(0< c_J\le 1\) small enough
for the derivative saving, and put

\[
H_t=\max(1,\left\lfloor b^{c_J \min(t,r)/(8\cdot 2^{D_r+L_1})}\right\rfloor)
\]

Then \(H_t\le b^{r/8}\), so \textbf{only for \(1\le h\le H_t\)} one has
\(\Lambda_h\le b^{-7r}\), \(\lambda_h m_{\mathrm{full}}^k\ge b^{t/2}\),
and the same variation ratio as in (5.2). The three terms of Appendix A
are bounded by \(b^{-c_J \min(t,r)/2^k}\) after decreasing \(c_J\) if
needed; the first ratio's polynomial in \(B\) factor is paid by
\(m_{\mathrm{full}}\asymp b^r\). Repeating the geometric offset
summation and exact all-prefix assembly in Section 2 gives
\(O_J(N B_r/(rD_r))\) for every fixed word.

For \(\theta=m\) positive integer, \(h_m(x)=x^m/(1-x)^m\) is analytic at
zero and has nonpolynomial rational derivative, so its Taylor
coefficients satisfy (W1)-(W2). The universal schedule from Section 2
applies without alteration, proving pointwise normality at these
isolated parameters. For \(\theta=0\), \(h_0=1\). Hence one schedule
works for the full punctured real axis; claimed compact uniformity is
restricted as stated. \(\square\)

\textbf{Important distinction.} For negative \(\theta\) the input
function diverges as \(x\to 0\), but its \textbf{fractional part} still
obeys the phase derivative argument; singularity is not a proof of
failure. At a positive integer, the fractional function itself is
rational and admitted, whereas its leading singular-power derivative
approximation vanishes and the separate analytic proof is needed. No
claim that mere nonpolynomiality of an arbitrary nonsmooth germ suffices
is made.

\section{Flat sources: suppression and
rescue}\label{flat-sources-suppression-and-rescue}

For \(g(x)=x/(1-x)\), the compositions

\[
\mathcal B_1(g(x))=\frac{g(x)}{1+g(x)}=x,\qquad \mathcal B_2(g(x))=\frac1{1+g(x)}=1-x
\]

are exact polynomial identities. Likewise \(\exp(\log(1+x))=1+x\). For
an integer-coefficient polynomial \(P\), the source values have at most
\(\deg(P)A_r+O(W)\) potentially nonconstant starts per copied row. All
these examples therefore fail in this observation model at the stated
fixed center for \textbf{every} admissible schedule with
\(D_r\to \infty\). Invertibility and a nonzero first derivative are
plainly insufficient.

\Needspace{6\baselineskip}
\textbf{Theorem G (flat, schedule-dependent survival).} Fix
\(g(x)=x/(1-x)>0\) for \(0< x< 1\) and
\(\mathcal B(\theta,y)=\exp(-\theta/y)\), with real \(\theta\ne 0\).
There exists one computable schedule independent of \(\theta\),
satisfying all source limits, for which every fixed nonzero \(\theta\)
produces a normal output. Convergence is uniform on each compact
interval avoiding zero. By contrast, under a \emph{balanced} schedule,
which has \(K_r=O(r \log r)\) and precision gap
\(B_r^{\mathrm{bal}}\ge c r/\log r\) for some fixed \(c>0\), the output
for each fixed \(\theta>0\) eventually consists entirely of zero copied
rows. Thus suppression by one schedule is not impossibility for the
operation.

\textbf{Proof.} The exact composite is
\(h_\theta(x)=\exp(\theta-\theta/x)\). Choose the following exact
integer schedule, for sufficiently large \(r\),

\[
B_r=3+\min\lbrace m\ge 1: b^{b^m}\ge r+4\rbrace ,\qquad A_r=r+B_r,
\]

\[
D_r=\max(1,\left\lfloor (\log_2 B_r)/8\right\rfloor).
\]

Use a bounded-increment smoothing of \(D\) and copy \(K_r=A_rD_r\)
digits. The threshold comparisons involve integers only, so this
schedule is computable without deciding equality of transcendental real
expressions. Then \(b^{B_r}=\Theta_b(\log r)\) and \(D_r\to \infty\).

Write the logarithm of the source as \(q(y)=\theta-\theta b^A/(by+1)\).
Its first derivative is \(q'(y)=\theta b^{A+1}/(by+1)^2\), of sign
\(\theta\), and its normalized higher derivatives satisfy

\[
\frac{|q^{(\ell)}(y)|}{|q'(y)|^\ell} \le  \ell! (C_b/(|\theta| b^{B_r}))^{\ell-1},\qquad \ell\ge 2,
\]

throughout the stage, with a fixed constant depending on \(b\).
Expanding \(\frac{d^k}{dy^k}\exp(q)\) by the elementary Bell partition
formula, the sum of terms containing an \(\ell\ge 2\) block has relative
magnitude at most

\[
\exp(C_b k^2/(|\theta|b^{B_r}))-1.\tag{6.1}
\]

For \(k\le D_r+12=O(\log B_r)\) this is below 1/2 eventually, uniformly
on \(\delta\le |\theta|\le M\). Hence its derivative of order \(k\) has
the sign of \(\theta^k\), and is comparable to \(\exp(q)|q'|^k\) by
factors between 1/2 and 3/2. Since
\(|q|=O_{\theta,b}(b^{B_r})=O_\theta(\log r)\) and
\(\log_b |q'|=-r+B_r+O_{\theta,b}(1)\), with stage variation bounded by
\(\exp(O_\theta(b^{B_r}))=r^{O_\theta(1)}\), the explicit choice
\(L=10\) in

\[
k=\left\lfloor t/(A_r-1)\right\rfloor+10
\]

supplies \(3\le k\le D_r+10\le D_r+12\) and, for
\(t\ge T_\theta b^{B_r}\), the harmonic-one bounds

\[
\lambda_1 m_{\mathrm{full}}^k \ge  b^{t/2},\qquad \Lambda_1\le b^{-8r},\qquad
\log_b(\frac{\Lambda_1}{\lambda_1})=O_\theta(\log r)
\]

Indeed \(t-k r\le -9r+O(D_r B_r)\) and the remaining logarithmic
derivative terms are \(O_\theta(D_r B_r+\log r)=o(r)\), leaving the
claimed \(-8r\) margin. Choose \(0< c\le 1\) small enough and

\[
H_t=\max(1,\left\lfloor b^{c \min(t,r)/(8\cdot 2^{D_r+12})}\right\rfloor)
\]

Since \(H_t\le b^{r/8}\), the upper bound \(\Lambda_h\le b^{-7r}\) is
asserted \textbf{only for \(1\le h\le H_t\)}; in that range the same
variation ratio and the lower bound
\(\lambda_h m_{\mathrm{full}}^k\ge b^{t/2}\) hold.

The order is at most \(D+12\), the growing-order estimate in Appendix A
applies, and its three terms are bounded by \(b^{-c \min(t,r)/2^k}\) for
a fixed positive \(c\) once \(r\) is large. Offset-dependent Fejér
cutoffs and summation as in Section 2, with
\(T_\theta b^{B_r}=O_\theta(\log r)\) early offsets, give the all-prefix
estimate

\[
|C_w(N;\theta)-Nb^{-W}| = O_{\theta,b,W}\left(\frac{N\log r(N)}{r(N)D_{r(N)}}\right).\tag{6.2}
\]

The constants and onset are uniform on compact subsets of
\(\mathbb R\setminus\lbrace 0\rbrace\). On that balanced schedule,
\(x_j\le b^{2-B_r^{\mathrm{bal}}}\), so
\(0< h_\theta(x_j)\le \exp(\theta-\theta b^{B_r^{\mathrm{bal}}-2})\) for
\(\theta>0\). Under the two comparison bounds stated above,
\(K_r=O(r \log r)\) and \(B_r^{\mathrm{bal}}\ge c r/\log r\), so
\(b^{B_r^{\mathrm{bal}}}/K_r\to \infty\). Hence this is eventually
smaller than \(b^{-K_r}\) for every row, forcing its full block of \(K\)
digits to zero. This suppression statement is restricted to that
schedule; it does not cover arbitrary copied lengths allowed to grow
faster than this exponential shrinkage. At \(\theta=0\) the source is
the integer 1. \(\square\)

\textbf{Reciprocal-square kernel, with the Bell estimate expanded.} For
\(\mathcal B(\theta,y)=\exp(-\theta/y^2)\), set

\[
z=by+1,\qquad a=b^A,\qquad v=a/z=1/x,\qquad q(y)=-\theta(v-1)^2
\]

Direct differentiation of \(-\theta(a^2 z^{-2}-2a z^{-1}+1)\) gives, for
every \(\ell\ge 1\),

\[
q^{(\ell)}=(-1)^{\ell+1} \theta \ell! b^\ell z^{-\ell} [(\ell+1)v^2-2v]
\]

In particular \(q'=2\theta(b/z)v(v-1)\) has sign \(\theta\) when
\(v\ge 2\). For \(\ell\ge 2\), using \(2v(v-1)\ge v^2\) and
\(\ell+1\le 4^{\ell-1}\),

\[
\frac{|q^{(\ell)}|}{|q'|^\ell} \le  \ell! [4/(|\theta|v^2)]^{\ell-1}
\le  \ell! [C_b/(|\theta|b^{2B_r})]^{\ell-1}
\]

To justify the summed Bell remainder, write
\(\eta=C_b/(|\theta|b^{2B_r})\). In the complete Bell partition formula
the leading term is \((q')^k\). Group all other terms by their number
\(j\) of blocks. The normalized absolute sum is bounded by

\[
\sum_{j=1}^{k-1} \frac{k!}{j!} \binom{k-1}{j-1} \eta^{k-j}
\le  \sum_{s=1}^{k-1} \frac{(k^2\eta)^s}{s!} \le  \exp(k^2 \eta)-1
\]

The coefficient in the first sum is the unsigned Lah number: it follows
by taking \(k!\) times the coefficient of degree \(k\) in
\(\exp(\sum_{\ell\ge 1} \eta^{\ell-1} u^\ell)\) in a formal variable
\(u\). Thus the relative remainder is at most

\[
\exp(C_b k^2/(|\theta|b^{2B_r}))-1,
\]

uniformly on compact parameter sets avoiding zero. This also justifies
(6.1) with \(\eta=C_b/(|\theta|b^{B_r})\) for the reciprocal kernel.

Choose

\[
B_r=3+\min\lbrace m\ge 1: b^{b^{2m}}\ge r+4\rbrace ,\qquad A_r=r+B_r,
\]

and the same bounded-increment depth with \(k\le D_r+12\). This
threshold again uses integer comparisons only. Then
\(b^{2B_r}=\Theta_b(\log r)\), the Bell remainder tends to zero,
\(|q|=O_\theta(\log r)\), and \(\log_b|q'|=-r+2B_r+O_\theta(1)\). The
same \(k=\left\lfloor t/(A_r-1)\right\rfloor+10\) gives
\(\Lambda_1\le b^{-8r}\) and
\(\lambda_1 m_{\mathrm{full}}^k\ge b^{t/2}\) after discarding
\(O_\theta(\log r)\) offsets; \(D_r B_r=o(r)\) pays the extra
\(2k B_r\). With the preceding \(H_t\), \(\Lambda_h\le b^{-7r}\) for
\(1\le h\le H_t\), and the variation logarithm is \(O_\theta(\log r)\).
The derivative estimate, offset sum and all-prefix assembly therefore
give (6.2) for this kernel as well. \textbf{This is a second flat-kernel
instance, not a blanket theorem for arbitrary \(C^\infty\) flat
operations.}

\section{The near-polynomial parameter
scale}\label{the-near-polynomial-parameter-scale}

Here \(\theta\) remains fixed along each output, but in a
\textbf{triangular comparison} we may choose a different parameter for a
different output cutoff \(N\); these are distinct infinite outputs, not
row-varying \(\theta_j\). Fix a nonnegative integer \(m\), a regular
stage schedule in Section 2 with \(K_{r-1}/K_r \to  1\), and
\(\alpha=m+\epsilon\) where \(0< |\epsilon|\le 1/2\). Let
\(E(\epsilon)=\max(0,\log_b(1/|\epsilon|))\). Offsets are zero-based:
offset \(t\) starts at digit \(t+1\), and a within-row word must satisfy
\(t\le K_r-W\). The sign of \(\epsilon\) is held fixed when taking a
sign-specific subsequence.

\Needspace{6\baselineskip}
\textbf{Proposition H (sharp scale for this schedule).} For any family
\(\epsilon_N\) with \(0< |\epsilon_N|\le 1/2\), the triangular word
frequencies of the outputs generated from
\(h_{m+\epsilon_N}(x)=(1+x)^{m+\epsilon_N}\) satisfy

\[
\frac{C_w(N;m+\epsilon_N)}{N} \to  b^{-|w|}\qquad\text{for every fixed }w
\]

\textbf{if and only if} \(E(\epsilon_N)/K_{r(N)} \to  0\). The analogous
statement holds in family (ii) with \(\theta=-m+\epsilon_N\), so the
actual exponent is \(m-\epsilon_N\) and the source is
\((1-x)^{m-\epsilon_N}\). In this convention positive \(\epsilon\) again
gives a positive perturbation and zero words, and negative \(\epsilon\)
gives borrowed \((b-1)\) words. The parameters are not allowed to vary
within one output.

\textbf{Sufficient direction.} For derivatives of order \(k>m\), the
exact identity

\[
\frac{d^k}{dx^k} (1+x)^{m+\epsilon} = (m+\epsilon)_{\underline{k}}(1+x)^{m+\epsilon-k}
\]

has magnitude bounded above and below, on \(0\le x\le b^{2-B_r}\), by
\(|\epsilon| \exp(\pm C_m k \log(k+2))\). For the derivative orders used
below, \(m+2< k\le D_r+L_m=O(\log B_r)\), the ratio of the maximum
magnitude to the minimum is at most \(\exp(C_m k b^{2-B_r})\) and is
therefore uniformly bounded at all sufficiently late stages, for
\(|\epsilon|\le1/2\). Choose explicitly \(L_m\ge \max(m+3,12)\) and
\(k=\left\lfloor t/(A_r-1)\right\rfloor+L_m\), so \textbf{every chosen
derivative order exceeds \(m\)}, and \(k\le D_r+L_m\) at all
sufficiently large stages. The whole high-derivative Taylor remainder
retains the \(\epsilon\) factor: define

\[
R_\epsilon(x)=\frac{(1+x)^{m+\epsilon}-(1+x)^m}{\epsilon}
\]

This has a uniformly bounded analytic derivative on a fixed interior
disk for \(|\epsilon|\le 1/2\), extending continuously at
\(\epsilon=0\). For every coefficient index \(n\ge m\),
\(c_n(\epsilon)=\epsilon a_n(\epsilon)\) with
\(|a_n(\epsilon)|\le M_m S_m^n\). Since \(k>m\), all coefficients in the
derivative remainder have that form, and (2.1) becomes

\[
\left|\frac{h^{(k)}(x)}{(k-1)!}-c_{k-1}(\epsilon)\right| \le  |\epsilon| 4 M_m k S_m^k x
\]

Thus the relative remainder estimate has no spurious factor
\(1/|\epsilon|\). Equivalently the exact derivative identity above
supplies the same uniform sign and variation bounds. The growing
derivative calculation in Section 2 works after discarding
\(O_m(B_r+E(\epsilon))\) actual offsets: take
\(t\ge T_m B_r+2 E(\epsilon)\) with a sufficiently large fixed \(T_m\).
The lower bound needed is the safe one

\[
\lambda_1 m_{\mathrm{full}}^k\ge b^{t/16}
\]

The harmonic-one upper bound is \(\Lambda_1\le b^{-8r}\); with an
offset-dependent cutoff as in Section 2 (and a sufficiently small fixed
saving constant), \(\Lambda_h\le b^{-7r}\) only for \(1\le h\le H_t\).
Fourier and harmonic summation pay \(O_m(B_r+E+2^D)\) starts per row and
the same ambient-row contribution. This bound is harmless even if the
retained range is empty: charge the row trivially, at most the
discarded-offset contribution. Summing with geometric row counts gives,
for all prefixes with \(r\) large and \(E=o(K_r)\),

\[
\frac{|C_w(N;m+\epsilon_N)-Nb^{-W}|}{N} \le  O_{m,b,W}\left(\frac{B_r}{rD_r}+\frac{E(\epsilon_N)}{K_r}\right)+o(1).\tag{7.1}
\]

For family (ii), replace the numerator defining \(R_\epsilon\) by
\((1-x)^{m-\epsilon}-(1-x)^m\); the same analytic \(\epsilon\) factor
and derivative argument apply.

This uses exact phases and Fejér intervals, so \textbf{no} numerical
digit-stability assumption is made near a carry boundary.

\textbf{Necessary direction and the borrow.} Let \(P_m(x)=(1+x)^m\) and
\(\Delta=(1+x)^{m+\epsilon}-P_m(x)\). For \(0< x\le 1/2\) and
\(|\epsilon|\le 1/2\),

\[
0< |\Delta|\le C_m |\epsilon| x,\qquad \mathrm{sgn} \Delta=\mathrm{sgn} \epsilon
\]

For row \(j\) in a late stage, \(P_m(x_j)\) is a rational polynomial
with denominator dividing \(b^{m A_r}\) and \(x_j\le b^{2-B_r}\). At
offsets \(t\ge m A_r\), \(b^t P_m(x_j)\) is an integer. Fix
\(C_{\ast}>2+\log_b C_m\). For any actual within-row offset satisfying

\[
m A_r\le t\le K_r-W,\qquad t+W\le E(\epsilon)+B_r-C_{\ast},
\]

one has \(b^t P_m(x_j)\) integer and the strict bound

\[
0< b^t|\Delta|\le C_m b^{t-E+2-B_r}< b^{-W}
\]

For \(\epsilon>0\) the fractional part of \(b^t(P_m+\Delta)\) lies in
the cylinder of \(W\) zeros. For \(\epsilon< 0\) it is \(1-b^t|\Delta|\)
and lies strictly inside the cylinder of \(W\) digits \(b-1\): this is
the \textbf{borrow}, with strict endpoint slack. This proves a digit
statement at a terminating rational endpoint, not a general inference
from closeness of two reals. The number of forced constant-word starts
per row is at least

\[
\max(0, \min(K_r,E+B_r-C_{\ast}-W)-m A_r-O(W)).\tag{7.2}
\]

If \(\limsup E(\epsilon_N)/K_{r(N)}>0\), pass to a sign-constant
subsequence with \(E\ge \rho K_r\). The regular schedule gives
\(N\le C_b b^r K_r\). Because \(D_r\to \infty\), \(m A_r=o(K_r)\),
\(B_r=o(K_r)\), \(K_{r-1}/K_r\to 1\), and the \textbf{entire previous
stage} is present at every such cutoff, (7.2) forces a positive fraction
\(\eta(\rho,b)>0\) of its starts to be the same constant word of any
fixed length \(W\). The preceding-stage proportion is bounded below
independently of each fixed \(W\); using \(K_{r-1}/K_r\to 1\), its
forced length is at least \(\min(1,\rho)K_{r-1}/2\) eventually. Choosing
one fixed \(W\) with \(b^{-W}< \eta/2\) contradicts normality. Thus
\(E/K_r\to 0\) is necessary. In family (ii),
\(\Delta=(1-x)^{m-\epsilon}-(1-x)^m\) also has sign \(\epsilon\) and
\(0< |\Delta|\le C_m|\epsilon|x\), so the same integer-endpoint and
borrow proof applies with this stated parameter convention. The proof
uses a \emph{long constant word}, not merely an excess of single zeros
that might in principle be compensated elsewhere. \(\square\)

For a \textbf{fixed} \(\epsilon\ne 0\), \(E\) is constant and the
criterion gives survival. For \(\epsilon_N\) exponentially small in
\(K\), a positive density of long constant blocks survives. This is a
scale threshold for these specified outputs, \textbf{not} a
classification of all imaginable schedules or a theorem about parameter
sequences \(\theta_j\) varying from row to row.

\section{Quantitative resources and local
normality}\label{quantitative-resources-and-local-normality}

This section records results for the same complete-prefix model as
Sections 1--7. Fix an integer base \(b\ge2\), positive integer schedules
\(A_j,D_j\), and

\[
x_j=(bj+1)b^{-A_j},\qquad K_j=A_jD_j.
\]

The output is the concatenation of the first \(K_j\) canonical
fractional base \(b\) digits of \(\theta+g(x_j)\). Write
\(S_J=\sum_{j\le J}K_j\), and let \(C_w(N;\theta)\) count occurrences of
a fixed word \(w\) of length \(W\) whose starts are below \(N\),
including starts whose words end later. Set \(E_w=C_w-Nb^{-W}\). All
center-uniform statements below concern \(\sup_\theta\), using one
schedule for every center. They do not assert normality of the center or
of an individual sampled value.

The results use several different schedules. In particular, minimizing
digits copied from each source and minimizing the number of admitted
sources are different objectives and are not asserted simultaneously.

\subsection{Exact grids, full-prefix rates, and signed logarithmic
mixtures}\label{exact-grids-full-prefix-rates-and-signed-logarithmic-mixtures}

For each of

\[
g_1(x)=\log(1+x),\quad g_2(x)=4\arctan x,\quad
g_3(x)=\log(1+x)+\log(1+x/2)+\log(1+x/3),
\]

define

\[
q_b=4\!\prod_{p\in\lbrace 2,3\rbrace ,\,p\nmid b}p,
\qquad \lambda_b=2\lceil\log_2q_b\rceil.
\]

On the dyadic stage \(2^{r-1}\le j< 2^r\), use

\[
\ell_r=\min\lbrace \ell:b^\ell\ge(b+1)2^r\rbrace ,\quad
A_r=(D_r+1)^2(\ell_r+r+4),\quad K_r=A_rD_r,
\]

where \(D_r\to\infty\), \(0\le D_{r+1}-D_r\le1\), and
\(D_r\le\max(1,\lfloor r/\lambda_b\rfloor)\). Then, at every output
cutoff,

\[
\sup_\theta\max_{|w|=W}|E_w(N;\theta)|
       =O_{b,W}(N/D_{r(N)}).
\tag{R1}
\]

The fastest allowed depth gives \(A_j\asymp(\log j)^3\),
\(K_j\asymp(\log j)^4\), and error \(O(N/\log N)\). A depth of order
\(\log r\) gives \(K_j\asymp\log j(\log\log j)^3\) and error
\(O(N/\log\log N)\).

Here is the finite arithmetic underlying this conclusion. If \(a_k\) are
the Taylor coefficients of one of the kernels, let \(H_d\) be their
reduced denominator lcm through degree \(d\), let \(Q_d\) be its factor
coprime to \(b\), and put

\[
P_d=\prod_{d/2< p\le d,\ p\nmid66b}p.
\]

Then \(Q_d\le q_b^d\), \(H_d/Q_d\mid b^{d+\lfloor\log_2d\rfloor}\), and
\(\sum_dP_d^{-1}< \infty\). At a common actual offset \(t\), with
constant \(A\), the truncated phase is the polynomial

\[
\sum_{k\le d}a_k(bj+1)^kb^{t-kA}.
\]

For each selected prime \(p\), only the coefficient of degree \(p\) has
a denominator divisible by \(p\), and its primary polynomial is a
nonzero affine function modulo \(p\). The prime 11 is genuinely excluded
for \(g_3\), since \(pa_p\equiv11/6\pmod p\). Conditional on the other
prime-power coordinates, CRT therefore gives an exactly uniform
translated \(P_d\)-point circle grid. Its interval-count error is at
most one, for either endpoint convention. Over any \(m\) consecutive
rows, the corresponding interval error is at most \(m/P_d+Q_d\). No
Euler alignment of \(A\) is required.

For a rational \(\eta>0\), choose \(\tau\) with \(2b^{-\tau}\le\eta\).
In a batch with numerator height at most \(b^\ell\), retain

\[
dA+\beta\le t\le(d+1)A-(d+1)\ell-\tau-W,
\quad 1\le d< D,\quad \beta=D+\lfloor\log_2D\rfloor.
\]

Writing \(L_d=\max(0,A-\beta-(d+1)\ell-\tau-W+1)\) and
\(G=\sum_{d< D}L_d\), the exact batch error is bounded by

\[
mU+V,\qquad
U=AD-G+2\eta G+\sum_{d< D}L_d/P_d,
\quad V=\sum_{d< D}L_dQ_d.
\tag{R2}
\]

All row-crossing words are already included in \(AD-G\). The Taylor
error is at most \(\eta\); expanded and contracted circle intervals pay
its effect without excluding digit-boundary centers. Summing these batch
bounds and charging the actually consumed part of the last partial row
proves a finite all-prefix certificate.

For the stated stages, \(U_r/K_r=O(1/D_r)\),
\(V_r\le A_rq_b^{D_r}/(q_b-1)\), \(K_r/K_{r-1}\le6\), and
\(N\asymp2^rK_r\). The degree cap makes \(q_b^{D_r}=O_b(2^{r/2})\).
Bounded depth increments and a geometric sum give
\(\sum_{s\le r}2^{s-r}/D_s=O(1/D_r)\). These bounds prove (R1) through
every cutoff, including incomplete batches and the final row.

The same method extends to every nonempty rational logarithmic mixture

\[
g(x)=\sum_{i=1}^s c_i\log(1+x/a_i),
\]

with distinct positive integers \(a_i\) and nonzero rational \(c_i\).
Vanishing initial moments are allowed. If \(r_0\) is the first index
with \(\sum_i c_i a_i^{-r_0}\ne0\), Vandermonde invertibility gives
\(r_0\le s\). Suitable fresh primes yield nonconstant primary
polynomials of degree at most \(2r_0-1\). Repeated Cauchy--Schwarz, with
zero differencing shifts explicitly paid, gives cancellation at every
unit multiplier. With the exact nonunit factor
\((\gcd(P_d,h)/P_d)^{2^{-(2r_0-1)}}\), the elementary prime-product
bound and an endpoint-safe Fejer inequality yield a summable full-period
interval discrepancy. Substitution in (R2) proves the same rate families
with input-dependent constants. This is a proof sketch of the extension;
its additional ingredient is the summable full-period discrepancy for
the nonlinear primary polynomials. Separate normality of the summands is
not used.

\subsection{A sharp qualitative window
scale}\label{a-sharp-qualitative-window-scale}

For a window beginning at \(M\) in dyadic stage \(r\), the preceding
schedule satisfies, for \(1\le L\le M\),

\[
\sup_\theta\frac{|E_w(M+L;\theta)-E_w(M;\theta)|}{L}
\le C_{b,W}\left(\frac1{D_r}+\frac{A_rq_b^{D_r}}L\right).
\tag{R3}
\]

Indeed complete-stage endpoints grow by a factor greater than two, so
this window meets at most two consecutive stages. The complete rows form
consecutive intervals, to which (R2) applies; at most two edge rows cost
\(2K_{r+1}\). Since \(A_{r+1}/A_r\le3\), \(D_{r+1}\le D_r+1\), and
\(D\le q_b^D\), the two period contributions and edge rows are absorbed
in the second term of (R3).

For every computable nondecreasing integer \(F(M)\to\infty\), one may
choose the depth so that every window of length

\[
L\ge F(M)\log(M+2)
\]

has uniformly correct fixed-word frequencies as \(M\to\infty\). To see
this, let \(d_r^*\) be the largest positive integer with
\(d^3q_b^d\le F(2^{\max(r-2,0)})\), using 1 if no such integer exists,
and take the minimum of \(d_r^*\), the preceding depth plus one, and the
usual degree cap. These caps are nondecreasing and divergent. Since
\(A_r=O_b(rD_r^2)\), the second term in (R3) is \(O_b(1/D_r)\) at the
specified threshold. For \(L>M\), subtract the two global prefix bounds
instead. Thus all longer windows are included, not only those contained
within one row batch.

Conversely no admissible schedule with \(x_j\to0\) and \(D_j\to\infty\)
can have this center-uniform conclusion for all windows with length
\(O(\log(M+2))\). These limits force \(K_j/\log j\to\infty\). For every
constant \(C\), infinitely many rows satisfy \(K_j>C\log(S_j+2)\):
otherwise \(S_j\le C_0+Cj\log(S_j+2)\), which first implies \(S_j< j^2\)
eventually and then \(K_j=O(\log j)\), a contradiction. At
\(M=S_{j-1}\), choosing \(\theta=-g(x_j)\) makes the entire row zero. A
window contained there has zero-digit error exactly \(1-1/b\). The
obstructing center may depend on the window, as permitted by the
supremum. This does not rule out smaller windows at a particular fixed
center.

\subsection{Prescribed large growth and source
admissions}\label{prescribed-large-growth-and-source-admissions}

For each of the three kernels, a positive integer demand \(h(j)\) can be
met by a center-uniformly normal schedule with \(K_j\ge h(j)\) at every
original index if and only if

\[
\log h(j)=o(j).
\tag{R4}
\]

The construction can in fact satisfy \(A_j\ge h(j)\). If \(h\) is
computable and satisfies the promise, the schedule is computable without
a supplied convergence modulus. The theorem does not decide whether an
input program is total or satisfies the promise.

For necessity, row zeroing gives the exact inequality

\[
(1-1/b)K_j\le \mathcal R_0(S_j)+\mathcal R_0(S_{j-1}),
\quad \mathcal R_0(N)=\sup_\theta|E_0(N;\theta)|.
\]

Uniform normality implies \(K_j/S_{j-1}\to0\). Hence
\(\log S_j=O(1)+\sum_{i=2}^j\log(1+K_i/S_{i-1})=o(j)\), and
\(\log K_j=o(j)\).

For sufficiency, put \(q_n=\max_{i\le2^{n+2}}\lceil\log_2h(i)\rceil\).
Let \(G\) linearly interpolate \(G(2^n)=q_n\), and set

\[
k_j=\lceil G(j)+\sqrt j\rceil,\qquad A_j=8\cdot2^{k_j}.
\]

Partition the row indices into maximal constant-height runs, of lengths
\(m_s\) and heights \(k_s\). Put \(m_0=D_0=1\), and on run \(s\) use

\[
D_s=\min\left\lbrace D_{s-1}+1,
\max\left(1,\left\lfloor\log_{576}\min(m_{s-1},m_s)\right\rfloor\right)\right\rbrace .
\]

All choices are finite exact operations; a run of height \(k\) ends by
row \(k^2\). The promise implies \(q_n=o(2^n)\), so the slopes of \(G\)
tend to zero. Therefore the height increments are eventually one and
\(m_s\to\infty\); consequently \(D_s\to\infty\). Also \(A_j\ge h(j)\),
all source arguments lie in the analytic domain, and tend to zero. The
inequalities

\[
Q_{D_s}\le24\sqrt{\min(m_{s-1},m_s)},\qquad
K_s/K_{s-1}\le4
\]

hold eventually. In (R2), the nonperiod relative error tends to zero,
since the last index in run \(s\) is at most \(k_s^2\), the depth is
\(O(\log k_s)\), and \(A_s=8\cdot2^{k_s}\). Completed-run period errors
have relative size \(O(Q_{D_s}/m_s)\to0\). At any partial current run,
its period and edge error divided by the length of the preceding
complete run is at most \(4(24\sqrt{m_{s-1}}+1)/m_{s-1}\to0\). Weighted
averaging proves normality at every original cutoff.

The same criterion applies to eventual demands and to demands on
\(A_j\). On a specified cofinal set \(J\), the condition is
\(\log h(j)=o(j)\) along \(J\); impose demand 1 elsewhere. If the
cofinal set may be chosen, the exact condition is
\(\liminf_j\log h(j)/j=0\). Choose increasing indices \(j_r\) with
\(h(j_r)^r\le2^{j_r}\), then use the preceding construction. This search
terminates under the promise. In contrast, the reindexed demand
\(K_{j_r}\ge h(r)\) is always achievable by passing to a sufficiently
sparse subsequence. Prescribing \(D_j\) itself remains a different,
unresolved classification.

Let \(J(N)=\min\lbrace j:S_j\ge N\rbrace\). Every center-uniform
normalizer satisfies

\[
J(N)/\log N\to\infty.
\tag{R5}
\]

This follows immediately from \(\log S_j=o(j)\). Conversely, for each
computable nondecreasing integer \(F(N)\to\infty\), the three kernels
admit a schedule with \(J(N)\le F(N)\log N\) eventually. Set
\(G(j)=\max(1,\min(F(j),\lfloor\sqrt j\rfloor))\) and demand
\(h(j)=2^{\lceil4j/G(j)\rceil}\). This is subexponential. For
\(S_j\le N< S_{j+1}\),

\[
F(N)\log N\ge F(j)\log h(j)
 \ge4(\log2)j\,F(j)/G(j)\ge j+1\ge J(N).
\]

These very long copied rows minimize admissions; they are not the
short-block construction below.

\subsection{Balanced near-logarithmic copied
length}\label{balanced-near-logarithmic-copied-length}

For every nonpolynomial rational derivative \(R\in\mathbb Q(x)\) regular
at zero, with primitive \(g(0)=0\), and every computable nondecreasing
integer \(F(j)\to\infty\), there is an effective schedule with

\[
A_j/\log_bj\to1,\qquad K_j\le F(j)\log_bj
\quad\text{eventually},
\]

\[
\sup_\theta\max_{|w|=W}|E_w(N;\theta)|
       =O_{R,b,W}(N/\log\log N).
\tag{R6}
\]

The implied constant is independent of \(F\). The argument uses the
growing-order derivative estimate in Appendix A. The mechanism is to
balance the derivative order against the precision gap.

\textbf{Proof outline.} The balancing mechanism is useful to state
explicitly. On the base \(b\) stage \(b^{r-1}\le j< b^r\), set
\(L_r=\max(1,\lfloor\log_2r\rfloor)\), choose nondecreasing divergent
depths with increments at most one and
\(D_r\le\max(1,\lfloor L_r/16\rfloor)\), and use

\[
a_r=\max(1,\lfloor L_r/(16D_r)\rfloor),\quad
B_r=c_0+\lceil r/a_r\rceil,\quad A_r=r+B_r.
\]

Here \(c_0\) is an explicit input-dependent analytic-radius guard. The
effective derivative order is at most \(\kappa_r=(a_r+1)D_r\le L_r/8\),
whereas \(a_rD_r\ge L_r/32\). Thus the discarded low offsets have
relative mass \(O(1/L_r)\), even when the copied depth diverges much
more slowly. Put \(R_r=b^{r-1}\) and \(M_r=(b-1)R_r\). The harmonic
cutoff is \(L_r^2\), and the actual-phase saving is
\(12M_rR_r^{-1/(a_r2^{\kappa_r})}\), uniformly on every consecutive row
subinterval. The coefficient recurrence gives an exponential upper
bound, a nonzero coefficient in each fixed-width window, and an
exponential lower bound for every nonzero coefficient. These control the
derivative sign and magnitude on the shrinking input interval. Applying
Appendix A at each retained offset, and Appendix B with cutoff
\(L_r^2\), gives the stated estimate after summing the discarded
offsets, the Fourier error, and the two edge rows.

Floors can make \(A_r\) decrease. The valid comparisons are
\(K_s\le2K_r\) for \(s\le r\) and \(K_r/K_{r-1}\le4\), not monotonicity
of precision. These suffice for every-prefix assembly. Adding the depth
caps \(\sqrt{L_r}\) and \(F(b^{r-1})/(c_0+4)\) forces \(a_r\to\infty\),
yielding \(A_r/r\to1\) and the prescribed copied-length bound without
changing the discrepancy constant.

The source limits alone imply \(K_j/\log j\to\infty\), so (R6)
approaches but does not attain a constant-overhead boundary. In the
class defined only by eventual limits and normality, there is no
coordinatewise minimal or lexicographically least schedule: decrease one
depth greater than one to one. Only a finite word is shortened; the
later sequence differs by a fixed shift, so all normality and limit
conditions survive. This observation does not apply after imposing extra
monotonicity, a fixed recurrence, or a finite accuracy constraint.

\subsection{Matching rate lower bounds and a genuine failure
regime}\label{matching-rate-lower-bounds-and-a-genuine-failure-regime}

The following exact center-averaging lemma prevents the early zero
excess from being dismissed as something that later digits might cancel.
If a batch of \(m\) equal-length rows has copied length \(K\), source
values \(\theta+y_j\), and \(\max y_j-\min y_j\le b^{-s}\) with integer
\(2\le s\le K\), then

\[
\sup_\theta E_{0,\mathrm{batch}}(\theta)
 \ge m[(1-1/b)s-1].
\tag{R7}
\]

Average \(\theta\) over \([-\min y_j,-\min y_j+b^{-s}]\). The first
\(s-1\) digits of every row are zero almost everywhere. The next
centered digit contributes at worst \(-1/b\). At every deeper offset the
center phase traverses an integer number of periods, so its centered
mean is exactly zero. Canonical endpoint choices affect only a null set.
Summing proves (R7).

Let \(N_r=S_{b^r-1}\) be the cumulative endpoint of stage \(r\). Then
the exact one-digit identity implies
\(\mathcal R_0(N_r)+\mathcal R_0(N_{r-1})\ge\sup_\theta E_{0,\mathrm{batch}}\).
For balanced schedules, bounded \(g'\) permits \(s=B_r-O_R(1)\). Since
\(B_r/K_r\ge1/\kappa_r\ge8/L_r\) and adjacent cumulative lengths are
\(O_b(M_rK_r)\), at least one of each sufficiently late pair of
endpoints satisfies

\[
\mathcal R_0(N)/N\ge (b-1)^2/(b^2L_r).
\]

Thus (R6) has matching order \(N/\log\log N\) infinitely often within
this family. For the dyadic CRT family, \(|g'|\le4\) and
\(s=A_r-\ell_r-\lceil\log_b4\rceil\) give the analogous lower bound
\((1-1/b)/(8D_r)\) at one of the two endpoints. Consequently its
\(N/D_r\) upper rate, including the fast \(N/\log N\) rate, also has
matching order. These are center-supremum, infinitely-often,
family-specific bounds. They are not lower bounds at a fixed center, at
every cutoff, or over all schedules.

Admission of a function does not mean every growth profile works. For
\(g(x)=x/(1-x)\), which is admitted by (R6), suppose

\[
D_j\log_b(bj+1)/A_j\to0.
\tag{R8}
\]

At center zero its zero-digit frequency tends to one at completed-row
endpoints, so the output is not normal. Put \(n=bj+1\), \(A=A_j\),
\(D=D_j\). Eventually \(2n^{D+1}< b^A\), and the positive geometric tail
is less than one after multiplying by \(b^{AD}\). Hence

\[
\lfloor b^{AD}g(n/b^A)\rfloor
       =\sum_{m=1}^D n^m b^{(D-m)A}.
\]

Each \(n^m\) fits in one field of \(A\) digits, with no interfield
carry. The number of nonzero digits is at most \(D(D+1)\log_b n/2+D\),
whose fraction of \(AD\) tends to zero. Weighted averaging proves the
endpoint claim. In particular the CRT precision geometry
\(A\asymp\log j\,D^2\) fails for this admitted kernel. The
prescribed-growth/window theorems above must not be generalized to all
rational derivatives merely by citing admission.

\subsection{An exact finite optimum with a stopping
proof}\label{an-exact-finite-optimum-with-a-stopping-proof}

For one fixed complete dyadic batch, put \(m=2^{r-1}\), \(\ell=\ell_r\),
\(\eta=2^{-r}\), and choose the least \(\tau\) with
\(2b^{-\tau}\le\eta\), and set \(\beta_D=D+\lfloor\log_2D\rfloor\).
Define

\[
\Gamma_D=1+\sum_{d< D}P_d^{-1},\qquad
T_D=(D-1)(\beta_D+\tau+W)+(D-1)(D+2)\ell/2.
\]

The sufficient certificate is

\[
\Phi(A,D)=\Gamma_D/D+2\eta+
 \frac{\sum_{d< D}Q_d}{mD}+\frac{T_D}{AD},
\quad A\ge\ell+1.
\tag{R9}
\]

For target \(0< \epsilon< 1\), let \(\sigma_D\) be \(\epsilon\) minus
the first three terms. A depth is feasible exactly when \(\sigma_D>0\),
and its least precision is

\[
A_{\min}(D)=\max\lbrace \ell+1,\lceil T_D/(D\sigma_D)\rceil\rbrace .
\]

Zero slack cannot be attained: for \(D\ge2\), \(T_D>0\), while \(D=1\)
already has fixed contribution exceeding one. The period contribution
\(B_D=(\sum_{d< D}Q_d)/(mD)\) is strictly increasing, since

\[
B_{D+1}-B_D=
\frac{DQ_D-\sum_{d< D}Q_d}{mD(D+1)}>0.
\]

It diverges because \(Q_d\ge P_d\to\infty\). Thus stopping at the first
\(B_D+2\eta\ge\epsilon\) excludes every later depth and proves a genuine
minimum of \(AD\) over all positive integer depths in (R9).

For the binary three-pole kernel, stage \(r=128\), \(W=3\), and
\(\epsilon=1/2\), this optimum is

\[
A=8205,\qquad D=19,\qquad K=155895.
\]

There are 36 feasible depths, 12 through 47; depth 48 triggers the
stopping rule. The sample count is \(2^{127}\) sources and is part of
the certificate. This optimizes a sufficient finite certificate, not
actual discrepancy or all analytic methods, and does not justify
optimizing infinitely many stages without a separate all-prefix assembly
proof.

\subsection{Limits and open
questions}\label{limits-and-open-questions}

A classification of prescribed depth and joint optimality of precision,
copied length, and source count remain open. The results concern newly
assembled digit streams. Identifying an output with an unchanged
traditional constant requires a separate argument; ordinary convergence
of sampled values does not transfer normality. Appendix C gives a
sufficient shifted-orbit condition and several obstructions.

\appendix

\section{Growing-order derivative
estimate}\label{growing-order-derivative-estimate}

Write \(e(t)=\exp(2\pi i t)\).

Arias de Reyna
\cite[Theorem 11, equation (15), and Remark 14, equation (19)]{arias}
proves the estimate below.

In the following weaker formulation all three constants are replaced by
12. If \(k\ge3\) and \(m>k\), and a real \(C^k\) phase satisfies

\[
0< \lambda\le f^{(k)}(x)\le\Lambda,
\]

on an interval containing m consecutive integer sampling points and the
unit intervals immediately preceding them, then, putting \(v=2^k\),

\[
\left|\sum e(f(j))\right|
 \le12m\max\left\lbrace 
  (\Lambda/(\lambda m))^{2/v},\quad
  (\Lambda^2/\lambda)^{1/(v-2)},\quad
  (\lambda m^k)^{-2/v}\right\rbrace .                 \tag{V1}
\]

The constants in that result are all below 11, uniformly in k. The
uniformity in k is essential here. A fixed-order derivative theorem
would not justify the construction. A derivative of constant negative
sign is handled by replacing f by -f.

The absolute right side of (V1) is nondecreasing as a function of m: the
three powers of m are respectively \(1-2/v\), \(1\), and \(1-2k/v\), all
nonnegative for \(k\ge3\). Thus replacing m by a larger ambient M is
valid. Short intervals \(m\le k\) are handled separately by the trivial
bound. This resolves the incomplete-stage issue without a lower bound on
the sampled row count.

\section{Endpoint-safe Fourier-to-words
lemma}\label{endpoint-safe-fourier-to-words-lemma}

For any \(n\ge1\) circle points \(z_1,\ldots,z_n\), any interval \(I\),
and \(H\ge1\), write \(F_h=\sum_{i=1}^n e(hz_i)\). Then

\[
\left|\frac{\#\lbrace i:z_i\in I\rbrace }n-|I|\right|
 \le\frac2{\sqrt{H+1}}+\frac2\pi\sum_{h=1}^H\frac{|F_h|}{hn}.
\tag{B.1}
\]

This intentionally elementary version is sufficient; no sharp
Erdős--Turán theorem is imported.

To prove it, use the Fejér kernel

\[
\mathcal F_H(t)=\frac1{H+1}\left|\sum_{a=0}^He(at)\right|^2.
\]

It is nonnegative, has mass one, and has Fourier coefficients
\(1-|h|/(H+1)\) for \(|h|\le H\). Since \(\sin(\pi t)\ge2t\) for
\(0\le t\le1/2\), its mass \(\alpha\) outside circle distance \(\delta\)
obeys

\[
\alpha\le[2(H+1)\delta]^{-1}.
\]

Convolve the closed expansion of \(I\) by \(\delta\) with
\(\mathcal F_H\), then add \(\alpha\); convolve the open contraction,
then subtract \(\alpha\). These are pointwise majorant and minorant,
including all endpoint conventions. Their constant terms differ from
\(|I|\) by at most \(2\delta+\alpha\); every nonzero Fourier coefficient
is at most \(1/(\pi|h|)\). Integrating against the empirical measure and
pairing the positive and negative frequencies gives (B.1) after taking
\(\delta=1/[2\sqrt{H+1}]\).

\section{Transfer and obstruction
principles}\label{transfer-and-obstruction-principles}

Normality of a constructed output does not prove normality of its
center. Ordinary convergence of source values does not supply the
shifted orbit comparison required below. Likewise, function-theoretic
transcendence and transcendence of a number are distinct questions.

Throughout, \(b\ge2\) is an integer, \(e(t)=\exp(2\pi i t)\),
\(\mathbb T=\mathbb R/\mathbb Z\), and
\(\lVert t\rVert =\mathrm{dist}(t,\mathbb Z)\). All digit expansions use
the terminating convention at radix-rational endpoints.

\subsection{Average orbit distance}\label{average-orbit-distance}

\Needspace{6\baselineskip}
\textbf{Proposition C.1 (Mean orbit equivalence).} The set

\[
\mathcal H_b=\left\lbrace z\in\mathbb T:
\frac1N\sum_{n< N}\lVert b^nz\rVert \longrightarrow0\right\rbrace
\]

is an additive subgroup. If \(x-y\in\mathcal H_b\), every fixed integer
character has the same limiting averages along the two radix orbits
whenever either limit exists. In particular, \(x\) is normal if and only
if \(y\) is normal.

\textbf{Proof.} The circle triangle inequality gives closure under
addition, and \(\lVert -z\rVert =\lVert z\rVert\) gives closure under
inverses. Also

\[
|e(hb^nx)-e(hb^ny)|\le2\pi|h|\lVert b^n(x-y)\rVert .
\]

Averaging and applying the Weyl criterion proves the conclusion.
\(\square\)

This is a valid statistical equivalence relation. It does not make the
set of normal numbers a subgroup: for any normal \(x\), both \(x\) and
\(-x\) are normal while their sum is zero.

\Needspace{6\baselineskip}
\textbf{Proposition C.2 (Shifted shadowing).} Suppose a real \(F\) is
made by concatenating complete prefixes, of lengths \(K_j\), of source
numbers \(s_j\). Put \(S_j=\sum_{i\le j}K_i\). For a prescribed target
\(C\), define

\[
\delta_j=\lVert b^{S_{j-1}}C-s_j\rVert ,
\qquad \ell_j=-\log_b\delta_j
\]

with \(\ell_j=+\infty\) if \(\delta_j=0\). Assume the number of rows is
\(o(N)\), the terminal row is negligible \(K_j/S_{j-1}\to0\), and

\[
\sum_{j\le J}(K_j-\ell_j)_+=o(S_J).
\tag{C.1}
\]

Then \(C-F\in\mathcal H_b\). Hence normality of \(F\) transfers to
\(C\).

\textbf{Proof.} The output and its row source have the same first
\(K_j\) digits at the row start. At local time \(0\le r< K_j\), their
circle distance is at most \(b^{r-K_j}\); summing this bound gives at
most \(1/(b-1)\) for that row. For the target and source, multiplication
by \(b\) is Lipschitz with constant \(b\) on the circle, so the distance
at time \(r\) is at most \(\min(1/2,b^r\delta_j)\). A geometric series
up to \(r\approx\ell_j\), followed by the trivial bound, yields

\[
\sum_{r< K_j}\min(1/2,b^r\delta_j)
\le C_b+(K_j-\ell_j)_+.
\]

The row triangle inequality, (C.1), and the \(o(N)\) row count prove
zero mean distance at row endpoints. An incomplete terminal row costs at
most \(K_j/2=o(N)\), proving the same for every prefix. \(\square\)

The start \(b^{S_{j-1}}C\) is essential. Replacing it by the unshifted
\(C\) changes the assertion. Ordinary numerical convergence \(s_j\to C\)
gives no such shifted shadowing estimate.

\subsection{An exact arithmetic grid for logarithmic
sources}\label{an-exact-arithmetic-grid-for-logarithmic-sources}

The following calculation explains the arithmetic mechanism for
logarithmic sources without imposing normality on the center.

\Needspace{6\baselineskip}
\textbf{Proposition C.3 (Logarithmic grids).} Fix integers
\(D\ge d\ge2\), let \(L_d=\mathrm{lcm}(1,\ldots,d)\), and let \(A\) be a
positive multiple of \(L_D\). Put

\[
n_j=bj+1,\quad x_j=n_jb^{-A},\quad
F_d(z)=\sum_{k=1}^d\frac{(-1)^{k-1}z^k}{k}.
\]

Let \(B\) be a nonnegative integer satisfying \(k_{(b)}\mid b^B\) for
all \(1\le k\le d\), where \(k_{(b)}\) is the part of \(k\) supported on
primes dividing \(b\). For \(t\ge dA+B\),

\[
b^tF_d(n_jb^{-A})\equiv b^tF_d(n_j)\pmod1.
\tag{C.2}
\]

For \(0< x_j\le1/2\) and \(t\le(d+1)A-E\),

\[
\left|b^t\log(1+x_j)-b^tF_d(x_j)\right|
\le2n_j^{d+1}b^{-E}.
\tag{C.3}
\]

Set

\[
P_d=\prod_{\substack{d/2< p\le d\\p\nmid b}}p.
\]

For any circle interval \(I\), any real \(\vartheta\), any integer
\(t\ge0\), and any consecutive interval \(T\) of integer row indices,

\[
\left|\#\lbrace j\in T:\lbrace \vartheta+b^tF_d(n_j)\rbrace \in I\rbrace 
-|T|\,|I|\right|
\le \frac{2|T|}{P_d}+2L_d.
\tag{C.4}
\]

The empty product is one. Equations (C.2) and (C.3), followed by
interval enlargement by the right side of (C.3), give a finite estimate
for the literal logarithmic phases.

\textbf{Proof.} Write \(k=k_{(b)}k_c\), where \((k_c,b)=1\). Since
\(\varphi(k_c)\le k\le D\), its value divides \(L_D\), so
\(b^A\equiv1\pmod{k_c}\). Both \(b^{t-Ak}\) and \(b^t\) vanish modulo
\(k_{(b)}\), and they agree modulo \(k_c\). Their difference is
therefore divisible by \(k\). Multiplying by \(n_j^k/k\) and summing
proves (C.2). Taylor's convergent series gives remainder at most
\(2x_j^{d+1}\); multiplication by \(b^t\) gives (C.3).

Write \(L_d=P_dQ\). Each prime in \(P_d\) occurs to the first power in
\(L_d\), since \(p>d/2\) and \(p^2>d\). In the integer polynomial
\(L_dF_d(z)\) modulo such a \(p\), every summand vanishes except
\(k=p\). Consequently it is a nonzero scalar multiple of
\(z^p\equiv z\pmod p\). Since \(p\nmid b\), the map \(j\mapsto n_j\) is
a permutation modulo \(p\). Conditional on a residue \(j\pmod Q\), all
the primary phase coordinates for \(p\) are independent uniform residues
by CRT. Their sum is a translated uniform grid of \(P_d\) points on the
circle. An interval contains \(|I|P_d+O(2)\) of those points. Thus (C.4)
holds over each complete period of length \(L_d\). The at most two
incomplete pieces of a consecutive row interval have total length below
\(2L_d\), proving the stated bound. \(\square\)

No statement about the digits of \(\vartheta\) entered this proof. All
denominator parts and the base factors were retained. An irregular row
subset would need its own discrepancy estimate modulo \(L_d\); (C.4)
does not automatically apply to it.

For a word of length \(\ell\), use the radix interval of length
\(b^{-\ell}\). A phase perturbation of size \(\eta\) changes membership
only within \(\eta\) of its two endpoints. Applying (C.4) to those
boundary intervals pays an additional \(O(\eta |T|+|T|/P_d+L_d)\). This
is a word-frequency argument with explicit boundary cost, rather than a
claim that approximate real values always share digits.

\subsection{Approximation and source count
barriers}\label{approximation-and-source-count-barriers}

\Needspace{6\baselineskip}
\textbf{Proposition C.4 (The complete-prefix row barrier).} Consider any
uniform-center whole-prefix construction

\[
X(\theta)=0.\mathrm{pref}_{b,K_1}(\theta+\delta_1)
\mathrm{pref}_{b,K_2}(\theta+\delta_2)\cdots
\]

with a schedule independent of \(\theta\). If fixed-word discrepancy
tends to zero uniformly in \(\theta\) through all prefixes, then
\(K_J/S_J\to0\). Consequently \(\log S_J=o(J)\) and
\(J/\log S_J\to\infty\).

\textbf{Proof.} Choose \(\theta_J=-\delta_J\pmod1\); the entire row
\(J\) consists of zeros. Uniform one-digit convergence at both \(S_J\)
and \(S_{J-1}\) gives

\[
(1-1/b)K_J=o(S_J)+o(S_{J-1})=o(S_J).
\]

This proves the first assertion. For \(J\ge2\),

\[
\log S_J-\log S_{J-1}=-\log(1-K_J/S_J)\longrightarrow0.
\]

Summing these increments and using Cesàro averaging yields
\(\log S_J=o(J)\). Since \(S_J\to\infty\), the last conclusion follows.
\(\square\)

This constrains the number of rows in this particular uniform
whole-prefix architecture. It is not a lower bound on the number of
distinct sources in all normal-number constructions, and says nothing
against a fixed normal number having a short independent definition.

\Needspace{6\baselineskip}
\textbf{Proposition C.5 (The initial-zero barrier).} Suppose
\(0< \epsilon_j< 1\), \(\epsilon_j\to0\), and \(X\) concatenates the
complete prefixes of \(K_j\) digits of \(\epsilon_j\) at center zero.
Set \(L_j=\max(0,\lceil\log_b(1/\epsilon_j)\rceil-1)\). If the row count
is \(o(S_J)\) and \(X\) is normal, then

\[
\sum_{j\le J}\min(K_j,L_j)=o(S_J).
\tag{C.5}
\]

\textbf{Proof.} The first \(L_j\) source digits are zero. Fix a word
length \(\ell\). There are at least
\(\sum_{j\le J}(\min(K_j,L_j)-\ell+1)_+\) starts of \(0^\ell\) in these
initial zero segments. If the left side of (C.5), divided by \(S_J\),
has limsup \(\rho>0\), these starts have limsup at least \(\rho\),
because \(\ell J=o(S_J)\). Normality gives the limit \(b^{-\ell}\).
Taking fixed \(\ell\) arbitrarily large forces \(\rho=0\). \(\square\)

For the logarithmic schedule, \(\epsilon_j\asymp j b^{-A_j}\), so the
closeness horizon is \(A_j-O(\log j)\), while the copied length is
\(A_jD_j\). Since \(D_j\to\infty\), most copied digits occur after this
horizon. Uniform normality of the output is therefore compatible with a
completely nonnormal center.

\Needspace{6\baselineskip}
\textbf{Proposition C.6 (Analytic rigidity).} Let \(g\) and \(h\) be
distinct real analytic germs at zero, with \(g(0)=h(0)=0\). Let
\(x_j=(bj+1)b^{-A_j}\to0\), and suppose \(A_j\to\infty\),
\(\log j=o(A_j)\), and \(K_j/A_j\to\infty\). Then \(g(x_j)\) and
\(h(x_j)\) cannot have identical first \(K_j\) fractional digits in base
\(b\) for infinitely many sufficiently large \(j\).

\textbf{Proof.} Their difference has a first nonzero Taylor term
\(cx^m\). Hence \(|g(x_j)-h(x_j)|\ge c'|x_j|^m\) for large \(j\).
Identical prefixes give circle distance at most \(b^{-K_j}\). Both
values tend to zero, so for large \(j\) their small real difference
equals their circle distance. But

\[
-\log_b(c'|x_j|^m)=mA_j-O(\log j)=O(A_j)=o(K_j),
\]

a contradiction. \(\square\)

Analytic interpolation cannot freely prescribe infinitely many long
blocks at these inputs. A proposed representation of a traditional
constant must supply the actual Taylor coefficients, their arithmetic
constraints, convergence, and the exponentially fine residual bounds.

\subsection{Local inverse approximation and arithmetic
separation}\label{local-inverse-approximation-and-arithmetic-separation}

\Needspace{6\baselineskip}
\textbf{Proposition C.7 (Local inverse approximation).} Let \(g\) be a
real \(C^1\) function on an interval, with \(g'(x_0)\ne0\), and let
\(\xi=g(x_0)\). On a sufficiently small compact neighborhood of \(x_0\),
superapproximation

\[
0< |\xi-g(p_n/q_n)|< H(p_n/q_n)^{-M_n},\quad
H\to\infty,\quad M_n\to\infty
\]

by reduced rationals of height \(H(p/q)=\max(|p|,q)\), with \(q>0\), is
equivalent to \(x_0\) being Liouville. Here all inputs must remain in
the regular inverse branch and converge to \(x_0\).

\textbf{Proof.} Continuity of the derivative and its nonvanishing give
constants \(0< c< C\) with

\[
c|r-x_0|\le |g(r)-g(x_0)|\le C|r-x_0|
\]

on that neighborhood. Since it is bounded, \(q\le H(p/q)\le C' q\) for
rational inputs in it. The fixed constants are absorbed by decreasing
each unbounded approximation exponent by a fixed factor. The positivity
of the error excludes \(r=x_0\). These observations prove both
implications. \(\square\)

For \(g=\log\) on positive inputs, this says that \(\xi\) is
superapproximable by rational logarithms exactly when \(e^\xi\) is
Liouville. At a fixed rational \(r=a/c>0\), distinct reduced \(p/q\)
satisfy

\[
|r-p/q|\ge1/(cq).
\]

The local lower Lipschitz bound for the logarithm therefore forbids such
superapproximation to \(\log r\). A construction requiring it cannot
identify its output with that prescribed logarithm.

\subsection{Classical normal Liouville
construction}\label{classical-normal-liouville-construction}

\subsubsection{Finite normal packets}\label{finite-normal-packets}

\Needspace{6\baselineskip}
\textbf{Proposition C.8 (Finite normal packets).} For each \(j\ge1\),
there is an effectively specified cyclic word \(W_j\) of length
\(L_j=b^j\) in which every word of length \(j\) appears exactly once
cyclically. Every word of length \(\ell\le j\) appears \(b^{j-\ell}\)
times cyclically. In the ordinary, noncyclic word its count differs by
at most \(\ell-1\).

\textbf{Proof.} Form the directed graph whose vertices are words of
length \(j-1\), including the single empty word when \(j=1\). Each word
of length \(j\) is an edge from its first \(j-1\) letters to its last
\(j-1\) letters. Each vertex has indegree and outdegree \(b\). The graph
is strongly connected: append the letters of a desired vertex in order.
A finite balanced connected directed graph has an Euler circuit. One
proof follows unused outgoing edges until a circuit closes, then splices
in a new circuit at any visited vertex with unused edges; balance
prevents getting stuck elsewhere, and finiteness terminates the process.
Choosing starting vertex and edges in lexicographic order makes the
construction effective.

Read one letter per traversed edge around the circuit. Each word of
length \(j\) is one edge, proving the first statement. A fixed word of
length \(\ell\) has \(b^{j-\ell}\) extensions to a word of length \(j\),
proving the second. Only the last \(\ell-1\) cyclic starts cross the end
of the ordinary word. \(\square\)

These classical de Bruijn packets provide explicit finite words for the
following argument.

\Needspace{6\baselineskip}
\textbf{Proposition C.9 (Repetition criterion).} Let \(W_j\) be words of
lengths \(L_j\to\infty\) such that the cyclic frequency of every fixed
word in \(W_j\) tends to its uniform value. Choose positive repetition
counts \(R_j\), and set

\[
U_j=\sum_{i< j}R_iL_i,
\qquad Y=0.W_1^{R_1}W_2^{R_2}\cdots.
\]

If \(L_j/U_j\to0\), then \(Y\) is normal through every prefix. The
finitely many initial stages with \(U_j=0\) are irrelevant.

\textbf{Proof.} Fix a word of length \(\ell\). The count in \(q\)
consecutive copies of \(W_j\) differs from its expected value by at most

\[
qL_j\epsilon_j+O(\ell),
\]

where \(\epsilon_j\to0\) is the cyclic frequency error; this follows by
considering the periodic sequence \(W_j^\infty\) and deleting its last
\(\ell-1\) starts. At a prefix inside stage \(j\), at most one
unfinished copy costs \(L_j\) starts. Stage joins contribute at most
\(\ell\) per stage. Because \(L_i\to\infty\) and \(R_i\ge1\), the number
of stages divided by the total length tends to zero. The weighted sum of
\(\epsilon_i\) over complete copies is \(o(N)\): discard finitely many
initial stages, then bound every later \(\epsilon_i\) by a common
arbitrary tolerance. The unfinished copy contributes at most
\(L_j/N\le L_j/U_j\to0\). These estimates hold at every prefix \(N\),
including the first few copies of a new packet. \(\square\)

\Needspace{6\baselineskip}
\textbf{Proposition C.10 (Normal Liouville numbers).} There is an
explicitly computable Liouville number normal in base \(b\).

\textbf{Proof.} Use the packets of Proposition C.8. Having computed
their lengths, choose \(R_j\) effectively large enough that, with
\(T_j=U_j+R_jL_j\),

\[
T_j>(j+1)(U_j+L_j),\qquad T_j\ge j L_{j+1}.
\]

Then \(L_{j+1}/U_{j+1}\le1/j\), so Proposition C.9 gives normality of
\(Y\). Let \(y_j\) be the rational with the first \(U_j\) digits of
\(Y\) and then the packet \(W_j\) repeated forever. It has a denominator
dividing

\[
Q_j=b^{U_j}(b^{L_j}-1)< b^{U_j+L_j}.
\]

Since \(Y\) and \(y_j\) agree for \(T_j\) digits,

\[
|Y-y_j|\le b^{-T_j}< Q_j^{-(j+1)}.
\]

Normality implies irrationality: an eventually periodic digit sequence
omits some sufficiently long word, since a period of length \(P\) has at
most \(P\) distinct starts of any given length in its tail. Thus the
displayed errors are nonzero. Reducing the denominators only strengthens
the approximation; the reduced denominators must be unbounded, since
otherwise a subsequence of these rational approximants would be constant
while converging to the irrational \(Y\). Therefore \(Y\) is Liouville.
\(\square\)

The factors \(b^{L_j}-1\) allow very long repetitions of packets which
already have the right finite statistics. Approximants with denominator
only \(b^n\) instead force long constant runs. This distinction is why
an excellent rational approximation need not obstruct normality.

\begin{thebibliography}{9}
\addcontentsline{toc}{section}{References}
\bibitem{arias} Juan Arias de Reyna, \emph{Explicit van der Corput's d-th derivative estimate}, arXiv:2407.02094v1 (2024), Theorem 11 and Remark 14. \url{https://arxiv.org/abs/2407.02094v1}.
\bibitem{becher} Ver\'onica Becher, Pablo Ariel Heiber and Theodore A. Slaman, \emph{A computable absolutely normal Liouville number}, manuscript, 30 January 2014. \url{https://math.berkeley.edu/~slaman/papers/liouville.pdf}.
\bibitem{madritsch} Manfred G. Madritsch, J\"org M. Thuswaldner and Robert F. Tichy, \emph{Normality of numbers generated by the values of entire functions}, Journal of Number Theory \textbf{128} (2008), 1127--1145. \url{https://madritsch.perso.math.cnrs.fr/downloads/papers/entireFinal.pdf}.
\end{thebibliography}
\end{document}
